% =============================================================================
% data_methods.tex
% Data Gathering Methodology and Diagnostics for Nobel Prize Multilayer Network
% Include in manuscript with: % =============================================================================
% data_methods.tex
% Data Gathering Methodology and Diagnostics for Nobel Prize Multilayer Network
% Include in manuscript with: % =============================================================================
% data_methods.tex
% Data Gathering Methodology and Diagnostics for Nobel Prize Multilayer Network
% Include in manuscript with: % =============================================================================
% data_methods.tex
% Data Gathering Methodology and Diagnostics for Nobel Prize Multilayer Network
% Include in manuscript with: \input{data_methods}
% =============================================================================

\section{Data Collection and Processing}\label{sec:data}

We constructed a multilayer directed network representing the institutional
structure of the Nobel Prize selection process across all five original prize
categories: Physics, Chemistry, Physiology or Medicine, Literature, and Peace.
The network comprises four layers---governing bodies, vetting bodies
(Nobel Committees), nominators, and nominees/laureates---connected by directed
edges representing institutional authority and participation. All data
gathering was performed programmatically in R, using a reproducible pipeline
of eight scripts. Source code is available at [repository URL]. Below we
describe each stage of the pipeline in sufficient detail for independent
replication.

\subsection{Governing Bodies}\label{sec:governing}

Governing bodies are the institutions that formally award each Nobel Prize.
We gathered membership rosters for four bodies, each mapped to its
corresponding prize category.

\subsubsection{Royal Swedish Academy of Sciences (RSAS)}

The RSAS awards the prizes in Physics and Chemistry. Members are elected for
life; their service begins at induction and ends at death. We scraped the
membership list from the Swedish-language Wikipedia page
\emph{Lista \"over ledam\"oter av Kungliga Vetenskapsakademien}, which
organizes members chronologically under bold year headers (e.g.,
\texttt{<p><b>1740</b></p>}) followed by bulleted lists of member names
linked to individual Wikipedia biography pages. Founding members (1739)
appear under the headings ``Akademiens grundare'' and ``\"Ovriga.''

We parsed the rendered HTML by iterating through child elements of the
main content division in document order, updating the current induction year
whenever a bold year header was encountered, and extracting all hyperlinked
names from subsequent unordered lists. We excluded redlinks (articles not
yet created) and non-article links. For each member, we recorded the
induction year as the start year; end years were left initially blank and
later imputed from Wikidata death dates (Section~\ref{sec:demographics}).

Wikidata QIDs were resolved in batch using the MediaWiki
\texttt{action=query\&prop=pageprops} API, which maps up to 50 Wikipedia
page titles to their corresponding Wikidata items per request. This is
approximately 50 times faster than scraping individual sidebar links. The
API also handles title normalization (e.g., converting underscores to spaces
and resolving Unicode variants), which we accounted for when matching
responses back to input titles. We used exponential backoff (delay doubling
from 2 seconds, up to 5 retries) for failed requests and imposed a
0.2-second delay between batches.

Where multiple records existed for the same QID within the RSAS (e.g., due
to class transfers within the Academy), we retained only the earliest
induction record. The final dataset contains 2{,}312 RSAS members spanning
1739--2025.

\subsubsection{Swedish Academy}

The Swedish Academy awards the Prize in Literature. Members are appointed for
life. We queried Wikidata directly via SPARQL, selecting all individuals
holding position (P39) ``member of the Swedish Academy'' (Q207360), excluding
honorary members (Q97563901). Start and end dates were extracted from the
\texttt{pq:P580} and \texttt{pq:P582} qualifiers on each position statement.
Years were parsed from Wikidata's ISO~8601 datetime format
(e.g., \texttt{1871-12-20T00:00:00Z} $\to$ 1871). The query yielded 202
member records.

All SPARQL queries throughout the pipeline were executed via HTTP POST to the
Wikidata Query Service endpoint (\texttt{https://query.wikidata.org/sparql}),
with results returned in \texttt{application/sparql-results+json} format. We
implemented retry logic with exponential backoff (5-second base delay, up to
10 retries) and identified our bot with a descriptive User-Agent string per
Wikidata's access policy.

\subsubsection{Storting (Norwegian Parliament)}

The Storting awards the Peace Prize. We queried Wikidata for all entities
classified as human (P31 = Q5) holding position (P39) ``member of the
Storting'' (Q9045502), extracting start and end dates from position
qualifiers. Norwegian parliamentarians may serve non-consecutive terms;
we consolidated consecutive terms by grouping each member's service
periods chronologically and merging periods where the gap between one
term's end year and the next term's start year was at most one year.
Non-consecutive terms were preserved as separate records. This yielded
4{,}281 member-term records spanning 1814--2025.

\subsubsection{Karolinska Institutet}

The Karolinska Institutet (KI) awards the Prize in Physiology or Medicine.
Prior to 1977, all full professors at KI constituted the selection body;
in 1977 the Nobel Assembly was established as a distinct entity, and in 1984
its membership was fixed at 50. Neither Wikidata nor KI's website publishes
a complete roster of the Nobel Assembly.

We compiled KI professor data from Project Runeberg's digitized editions of
the Swedish state calendar (\emph{statskalendern}), covering 1881--1970.
The resulting Excel file contains one row per professor with columns for
Wikidata QID, name, and the first and last calendar edition in which each
professor appeared (serving as proxies for start and end years of service).
We retained the 140 records with confirmed QIDs and both start and end years.

\paragraph{Known gap.} Coverage ends at 1970. Post-1970 Nobel Assembly
membership is not publicly available: Wikidata contains essentially no
P39/P463 records for the Nobel Assembly (Q3375124), KI's website does not
publish a member roster, and Project Runeberg's digitized calendars stop
circa 1972. Extending this coverage would require direct contact with the
Nobel Office at KI.

\subsection{Vetting Bodies (Nobel Committees)}\label{sec:vetting}

Each prize has a dedicated Nobel Committee that reviews nominations and
prepares recommendations for the governing body. We gathered membership
data for all five committees.

\subsubsection{Chemistry, Physics, and Physiology/Medicine Committees}

These three committees are documented on Swedish-language Wikipedia. We
extracted membership lists using a general-purpose parser that fetches
wikitext (not rendered HTML) via the MediaWiki API
(\texttt{action=parse\&prop=wikitext}), splits the text into sections by
heading markers (\texttt{==}~and~\texttt{===}), and extracts bullet-list
entries from sections titled ``Tidigare ledam\"oter'' (former members),
``Nuvarande ledam\"oter'' (current members), and
``Sekreterare'' (secretaries).

Each entry follows one of several Swedish Wikipedia conventions:
\texttt{[[Name]], YYYY--YYYY} (former members with date range),
\texttt{[[Name]], invald YYYY} (current members, Swedish format), or
\texttt{[[Name]], YYYY--} (current members, open-ended). We parsed
Wikipedia article titles from double-bracket links (handling both
\texttt{[[Title|Display Name]]} and \texttt{[[Title]]} forms), extracted
start and end years via regex, and resolved QIDs through the batch
MediaWiki API.

The Swedish Wikipedia page for the Physiology/Medicine committee explicitly
states ``Listan \"ar ofullst\"andig'' (the list is incomplete); our data
should be understood in light of this limitation. Final counts: Chemistry 50
members, Physics 51 members, Physiology/Medicine 72 members (59 with QIDs).

\subsubsection{Literature Committee}

The Nobel Committee for Literature is drawn from the Swedish Academy.
No standalone membership list exists on Wikipedia. Instead, we extracted
committee service years from individual member biographies on
\texttt{svenskaakademien.se}.

We first queried Wikidata for all Swedish Academy members with their
seat numbers (extracted from position labels matching the pattern
``member of the Swedish Academy at seat $N$''). We then constructed
biography page URLs following the site's convention:
\texttt{/de-aderton/stol-nr-\{seat\}-\{name-slug\}}, where the name
slug is a lowercased, hyphen-separated, ASCII-transliterated version
of the member's name (e.g., \"a$\to$a, \"o$\to$o, \aa$\to$a).

The Swedish Academy website is built as a Nuxt.js single-page application.
Rather than executing JavaScript, we extracted biography text from the
\texttt{\_\_NUXT\_DATA\_\_} script tag embedded in each page's HTML source.
This payload contains HTML-encoded biography text; we decoded HTML entities,
stripped markup tags, and searched for mentions of ``Nobelkommitt\'en''
followed by year ranges. Two-digit years were expanded to four digits
based on context (e.g., ``1969--86'' $\to$ 1969--1986). We retained only
members whose biographies mentioned Nobel Committee service, yielding 35
members.

\subsubsection{Norwegian Nobel Committee}

The Norwegian Nobel Committee both vets nominations and awards the Peace
Prize. We parsed the membership table from the English Wikipedia page
\emph{List of members of the Norwegian Nobel Committee}. The page contains
a wikitable with columns for member name, start year, end year, tenure,
party affiliation, and chair/deputy status. We extracted the largest table,
parsed start and end years from the second and third columns, and resolved
QIDs from the first hyperlink in each row via the batch MediaWiki API.
This yielded 61 members (57 with QIDs).

\subsection{Nominations}\label{sec:nominations}

We scraped the Nobel Prize Nomination Archive
(\texttt{nobelprize.org/nomination/archive/}), which is subject to a 50-year
secrecy rule. As of 2025, nomination records are available through 1974 or
1975 for Physics, Chemistry, Literature, and Peace. Physiology or Medicine
nominations are available only through 1953; the Karolinska Institutet has
not released records for 1954--1974, likely related to the 1977 Swedish
transparency law and the concurrent restructuring that created the Nobel
Assembly.

The archive's web interface uses numeric prize codes in its URL parameters:
Physics = 1, Chemistry = 2, Physiology or Medicine = 3, Literature = 4,
Peace = 5. For each prize-year combination, we first retrieved the list page
(\texttt{list.php?prize=\{code\}\&year=\{year\}}) and extracted all
nomination IDs from links to detail pages (\texttt{show.php?id=\{id\}}).
We then scraped each detail page to extract the nomination year, prize
category, and the names and internal person IDs of all nominees and
nominators.

Detail pages use a two-column table layout with bold section headers
(``Nominee:'' and ``Nominator:'') followed by person links
(\texttt{show\_people.php?id=\{id\}}). We classified each person link
by walking through all table cells in document order and tracking which
section header was most recently encountered. A single nomination may list
multiple nominees (16.3\% of nominations) and, more rarely, multiple
nominators (1.1\%), reflecting joint nominations and co-signed nomination
letters. For multi-nominee or multi-nominator nominations, we generated
all pairwise nominee--nominator combinations.

Prize names were standardized to a common vocabulary (e.g., ``Physiology
or Medicine'' $\to$ ``Physiology/Medicine''; the Peace Prize archive uses
the header ``Nomination for Nobel Peace Prize'' rather than the
``Nobel Prize in\ldots'' format used by the other four categories, which we
handled as a special case in parsing).

List pages were fetched sequentially with 0.5-second inter-request delays;
detail pages were fetched in parallel across $N-1$ CPU cores (where $N$ is
the number of available cores, capped at 12) with 0.2-second per-worker
delays. Partial results were checkpointed to disk after each prize-year to
allow resumption if the process was interrupted. The final dataset contains
27{,}932 nomination records spanning 22{,}600 unique nominations, 4{,}006
unique nominees, and 11{,}650 unique nominators. Among these, 755 records
(2.7\%) have anonymous or unlisted nominators, and 2 records represent
self-nominations.

\subsection{Laureates}\label{sec:laureates}

We retrieved all Nobel Prize laureates from the official Nobel Prize API
(v2.1, \texttt{api.nobelprize.org/2.1/laureates}). The API provides
structured JSON data including Wikidata QIDs, names, gender, award year,
prize category, prize portion, motivation text, and primary affiliation at
the time of award. We paginated through all results (25 records per page)
and filtered to the five original prize categories, excluding the Sveriges
Riksbank Prize in Economic Sciences (established 1969).

Each laureate record was processed to extract the Wikidata QID (from the
\texttt{wikidata.id} field), full name (preferring \texttt{fullName.en}
over \texttt{knownName.en}), gender, and prize details. Laureates who
received multiple prizes (7 individuals, including Marie Curie, Linus
Pauling, Frederick Sanger, John Bardeen, and K.~Barry Sharpless) appear
as separate records per prize. The final dataset contains 927
laureate-prize records representing 919 unique individuals, with 100\%
QID coverage.

\subsection{Entity Resolution: Matching Nomination Archive People to
Wikidata}\label{sec:entity-resolution}

The nomination archive identifies people by internal numeric IDs and names
only, without Wikidata QIDs or other persistent identifiers. To integrate
nomination data with the rest of the network (which uses Wikidata QIDs as
the canonical identifier), we developed a two-stage matching pipeline.

\subsubsection{Stage 1: Biographical Data Extraction}

For each of the 14{,}471 unique people in the nomination archive, we
scraped their detail page on \texttt{nobelprize.org}
(\texttt{show\_people.php?id=\{id\}}), extracting last name, first name,
gender, birth year, and death year from labeled table cells. Scraping was
parallelized across available cores in chunks of 500, with partial results
checkpointed to disk.

Of the 14{,}471 people, 4{,}571 (31.6\%) had birth years recorded;
the remaining 9{,}903 (68.4\%) lacked birth year data entirely,
rendering them unmatchable through our disambiguation strategy.

Six records contained non-standard gender codes in the source database
(\texttt{<}, \texttt{\S}, \texttt{-}, \texttt{N}, \texttt{J}). We verified
the identities of all six individuals by first name and confirmed they are
male; these were recoded to ``M'' in our output. One record (Antonio
Carrillo Flores) had a transposed death year (1886 instead of 1986); this
was corrected manually.

\subsubsection{Stage 2: Wikidata Matching}

For each person with both a name and birth year (4{,}571 matchable
individuals), we searched the Wikidata API
(\texttt{wbsearchentities}, limited to 10 candidates per query) and then
verified each candidate's birth year via the \texttt{wbgetclaims} API
(property P569, date of birth). A match was accepted only when the
candidate's Wikidata birth year exactly equaled the birth year recorded
in the Nobel archive. We did not attempt fuzzy name matching or
approximate year matching, prioritizing precision over recall.

This procedure matched 339 of 4{,}571 matchable individuals (7.4\%),
representing 2.3\% of all 14{,}474 nomination archive people. The low
match rate reflects several factors: many early-twentieth-century
nominators and nominees lack Wikidata entries; name transliterations
between the archive's conventions and Wikidata's labels may differ; and
our strict birth-year criterion rejects plausible matches where Wikidata
records a different (or no) birth date. Unmatched individuals are
represented in the network using temporary identifiers with a \texttt{NOM:}
prefix (e.g., \texttt{NOM:7283}), preserving the nomination-layer topology
even where cross-layer entity resolution was not possible.

\subsection{Demographic Enrichment}\label{sec:demographics}

We fetched demographic attributes from Wikidata for all individuals with
resolved QIDs across all data sources. The SPARQL query retrieved: English
name label, gender (P21), country of birth (P19, traversed to sovereign
state via P17), nationality (P27), birth date (P569), death date (P570),
occupation (P106), and employer/institution (P108). Queries were executed
in batches of up to 200 QIDs using the \texttt{VALUES} clause, with
individual fallback queries for any batch that failed.

Multi-valued fields (nationality, occupation, institution) were collapsed
into semicolon-delimited strings. Where multiple birth or death dates
existed for a single QID (e.g., due to conflicting sources), we retained
the first non-missing value. Birth and death years were parsed from
Wikidata's ISO~8601 datetime format.

After the initial demographic fetch for governing body, vetting body, and
laureate QIDs, we performed a second pass to fetch demographics for the
339 nomination archive people matched in Section~\ref{sec:entity-resolution},
then re-collapsed all demographic records to produce a single row per QID.

As a post-processing step, we imputed end-of-service years for RSAS members
(who serve for life) using Wikidata death years. Of 2{,}312 RSAS members,
1{,}868 received imputed end years; 444 remain without end years (likely
living members or individuals with unknown death dates).

The final demographics table contains 7{,}221 unique individuals with
coverage rates of 99.8\% for name, 99.5\% for gender (90.3\% male, 9.3\%
female), 78.7\% for birth country, 97.2\% for nationality, 99.5\% for
birth year, 78.2\% for death year, 97.4\% for occupation, and 38.4\% for
institutional affiliation.

\subsection{Network Construction}\label{sec:network}

The final network was assembled from all intermediate data files. Nodes
correspond to unique individuals identified by Wikidata QIDs, with
demographic attributes from Section~\ref{sec:demographics}. Edges are
directed and carry metadata for year, prize category, and the layer
identities of the source and target nodes.

We defined five edge types reflecting the institutional flow of the
selection process:

\begin{enumerate}
    \item \textbf{Governing body $\to$ vetting body.} For each vetting
    body member, we created edges from all governing body members who were
    active at the start of the vetting member's service. ``Active'' means
    the governing body member's start year is at or before the vetting
    member's start year, and the governing body member's end year (if known)
    is at or after the vetting member's start year. The edge year is set to
    the vetting member's start year.

    \item \textbf{Vetting body $\to$ nominator.} For Chemistry, Physics,
    and Physiology/Medicine, Nobel Committees solicit nominations from
    qualified individuals. For each prize-year, we created edges from all
    active vetting body members to all nominators who submitted nominations
    that year. Nominators without resolved QIDs were represented using
    \texttt{NOM:} prefix identifiers.

    \item \textbf{Nominator $\to$ nominee.} Direct edges from each
    nominator to each nominee within a nomination record. For nominations
    with multiple nominees or nominators, all pairwise combinations are
    represented. Edges were deduplicated.

    \item \textbf{Governing body $\to$ laureate.} For Chemistry, Physics,
    Physiology/Medicine, and Literature, we created edges from all governing
    body members active in the award year to the laureate. This represents
    the formal decision-making authority of the governing body.

    \item \textbf{Vetting body $\to$ laureate.} For Peace only, the
    Norwegian Nobel Committee both evaluates nominations and awards the
    prize. Edges connect all active committee members in the award year
    to the laureate.
\end{enumerate}

The body-to-prize mappings are: RSAS $\to$ \{Chemistry, Physics\};
Karolinska Institutet $\to$ Physiology/Medicine; Swedish Academy $\to$
Literature; Storting $\to$ Peace. For vetting bodies: Nobel Committee for
Chemistry $\to$ Chemistry; Nobel Committee for Physics $\to$ Physics; Nobel
Committee for Physiology/Medicine $\to$ Physiology/Medicine; Nobel Committee
for Literature $\to$ Literature; Norwegian Nobel Committee $\to$ Peace.

After deduplication, the network contains 7{,}221 nodes and 305{,}199
edges. Edge counts by type: governing body $\to$ laureate (147{,}965;
48.5\%), vetting body $\to$ nominator (87{,}492; 28.7\%), governing body
$\to$ vetting body (42{,}360; 13.9\%), nominator $\to$ nominee (26{,}654;
8.7\%), and vetting body $\to$ laureate (728; 0.2\%). Edge counts by prize:
Physics (121{,}290), Chemistry (111{,}622), Physiology/Medicine (49{,}060),
Peace (15{,}853), Literature (7{,}374). The lower counts for
Physiology/Medicine reflect the Karolinska data gap post-1970 and the
truncated nomination archive (through 1953 only); the lower count for
Literature reflects the smaller Swedish Academy membership (18 seats).

\subsection{Self-Loops and Dual Roles}\label{sec:selfloops}

The network contains 220 self-loops (edges where the source and target
are the same individual). These arise legitimately from individuals who
occupy roles in multiple layers simultaneously. Specifically: 148
governing $\to$ vetting edges where RSAS members also served on Nobel
Committees; 44 governing $\to$ laureate edges where RSAS members
subsequently won Nobel Prizes (e.g., Svante Arrhenius, Ernest Rutherford);
26 vetting $\to$ nominator edges where committee members also submitted
nominations; and 2 nominator $\to$ nominee self-nominations. We retained
all self-loops as they reflect genuine structural features of the
selection process.

\section{Data Quality Assessment}\label{sec:diagnostics}

We performed a comprehensive 42-point diagnostic audit of all intermediate
and output data files. Below we summarize the results, organized by data
source, and assess the overall fitness of the dataset for network analysis.

\subsection{File Integrity and Completeness}

All nine expected data files (seven intermediate, two final output) were
present and parseable. Row counts, column counts, file sizes, and
modification timestamps were verified. No files were missing or corrupted.

\subsection{Governing Body Quality}

\paragraph{QID completeness.} All four governing bodies achieved 100\% QID
resolution: RSAS (2{,}312/2{,}312), Storting (4{,}281/4{,}281), Swedish
Academy (202/202), and Karolinska Institutet (140/140).

\paragraph{Year coverage.} RSAS membership spans 1739--2025; Storting spans
1814--2025; the Swedish Academy spans 1786--2024; Karolinska Institutet
spans 1884--1970. All year values fall within plausible historical ranges
(1700--2026). No records have start year exceeding end year (after
correction of three RSAS records where death-year imputation produced
anomalous end years; these were reverted to missing).

\paragraph{RSAS end-year imputation.} Of 2{,}312 RSAS members, 1{,}868
received end years imputed from Wikidata death dates. The remaining 444
(19.2\%) lack end years, likely because they are still living or their
death dates are not recorded in Wikidata. These members are treated as
currently active when determining temporal overlap for edge construction.

\paragraph{Multiple terms.} 583 QIDs appear multiple times within the same
governing body. For the Storting, these represent genuinely non-consecutive
parliamentary terms. For RSAS and Swedish Academy, only the earliest record
per QID was retained (duplicates arise from class transfers within the
Academy).

\subsection{Vetting Body Quality}

\paragraph{QID completeness.} Chemistry (50/50, 100\%), Physics (51/51,
100\%), Literature (35/35, 100\%), and the Norwegian Nobel Committee
(57/61, 93.4\%) all achieved high QID resolution. The Physiology/Medicine
committee achieved 59/72 (81.9\%), reflecting the acknowledged
incompleteness of the source list.

\paragraph{Cross-validation with governing bodies.} We verified that
committee members appear in their parent governing body's roster: 42/45
Chemistry committee members appear in RSAS (3 not found), 44/50 Physics
committee members appear in RSAS (6 not found), and 23/52
Physiology/Medicine committee members appear in Karolinska Institutet
(29 not found). The high rate of missing Physiology/Medicine cross-references
is attributable to the Karolinska data covering only 1881--1970, while the
committee list extends to 2025.

\paragraph{Known source limitations.} The Physiology/Medicine committee's
Swedish Wikipedia page states the list is incomplete. The Literature
committee data was extracted from the Swedish Academy's Nuxt.js-based
website by parsing the \texttt{\_\_NUXT\_DATA\_\_} payload, a method that
is fragile to future site redesigns.

\subsection{Nomination Data Quality}

\paragraph{Prize coverage.} All five prize categories are represented:
Physiology/Medicine (6{,}321 records), Peace (6{,}134), Physics (5{,}691),
Chemistry (5{,}094), and Literature (4{,}692). Year ranges align with
expectations: Chemistry, Physics, and Literature cover 1901--1974; Peace
covers 1901--1975; Physiology/Medicine covers 1901--1953.

\paragraph{Missing data.} All 27{,}932 records have nominee person IDs.
755 records (2.7\%) lack nominator person IDs, corresponding to anonymous
or unlisted nominators. Two self-nominations were detected (nominee =
nominator).

\paragraph{Entity resolution.} Of 14{,}474 unique people in the nomination
archive, only 339 (2.3\%) were matched to Wikidata QIDs. The primary
bottleneck is missing biographical data: 9{,}903 people (68.4\%) lack
birth years in the archive, making Wikidata disambiguation infeasible.
Of the 4{,}571 people with sufficient data for matching, 339 (7.4\%) were
successfully resolved. As a consequence, 21.4\% of all edge endpoints
in the final network use unresolved \texttt{NOM:} prefix identifiers, and
35.4\% of edges involve at least one such identifier. These edges preserve
the nomination-layer topology but cannot be linked to individuals in other
layers.

\subsection{Laureate Data Quality}

All 927 laureate-prize records have Wikidata QIDs (100\% coverage), and
all 919 unique laureate QIDs appear in the nodes file. Gender is recorded
for 891 laureates (827 male, 64 female, 28 unknown). Seven individuals
received multiple prizes. All laureates span 1901--2025.

\paragraph{Governing $\to$ laureate edge coverage.} Of 919 laureates, 267
lack governing body $\to$ laureate edges. Of these, 139 are Peace laureates
(expected, since Peace uses the vetting $\to$ laureate pathway), and 128
are Physiology/Medicine laureates awarded after 1970 (the Karolinska data
gap). These 128 missing edge sets correspond to 55 prize-years
(1971--2025) and represent a structural limitation of the available data
rather than a pipeline error.

\subsection{Demographic Data Quality}

The demographics table contains 7{,}221 unique QIDs with no duplicates.
All governing body QIDs (6{,}133), all vetting body QIDs (235), and all
matched nomination QIDs (339) are present. Completeness rates for individual
fields were reported in Section~\ref{sec:demographics}.

\paragraph{Temporal consistency.} No records have birth year exceeding
death year (after correction of one transposed death year). Fifty-seven
records have birth years prior to 1700, corresponding to early members
of the RSAS (founded 1739) and Storting (founded 1814); these are
historically accurate.

\paragraph{Gender distribution.} Demographics report 6{,}518 male (90.3\%),
670 female (9.3\%), and 33 unknown (0.5\%). The nomination archive's own
gender field reports 10{,}054 male (69.5\%), 334 female (2.3\%), and 4{,}086
unknown (28.2\%). The large unknown fraction in the nomination archive
reflects missing biographical data for people without Wikidata entries.

\subsection{Network Structural Quality}

\paragraph{Duplicate edges.} No duplicate edges exist after deduplication.

\paragraph{Isolate nodes.} 2{,}466 of 7{,}221 nodes (34.2\%) have no
edges. These are predominantly nomination archive people who were matched
to Wikidata QIDs (and thus appear in the nodes table) but whose
\texttt{NOM:} prefix identifiers in the edge list were not replaced by
their QIDs. This is a consequence of the network construction pipeline
resolving QIDs for nodes and edges independently; nomination-layer edges
reference the original \texttt{NOM:} identifiers regardless of whether
the person was subsequently matched.

\paragraph{Degree distribution.} Among edges with resolved QIDs on both
endpoints, 4{,}624 nodes are involved. The mean degree is 85.3, the median
is 13, the minimum is 1, and the maximum is 1{,}159 (QID Q4951154). The
high maximum degree reflects governing body members of large institutions
(e.g., RSAS members connected to many laureates and committee members over
long service periods).

\paragraph{Cross-layer entity overlap.} 105 laureates are also governing
body members, 15 laureates served on vetting bodies, 68 laureates appear
as nominees (with resolved QIDs), and 113 individuals appear in multiple
governing bodies. These overlaps reflect the historically intertwined nature
of Swedish scientific institutions and are expected features of the data.

\subsection{Summary of Known Structural Gaps}\label{sec:gaps}

We document seven structural limitations inherent to the available sources,
all of which should be considered when interpreting network analyses:

\begin{enumerate}
    \item \textbf{Karolinska Institutet post-1970 gap.} Nobel Assembly
    membership after 1970 is not publicly available, resulting in missing
    governing body $\to$ laureate and governing body $\to$ vetting body
    edges for Physiology/Medicine after 1970.

    \item \textbf{Nomination archive 50-year secrecy rule.} Nomination
    data are available through 1974--1975 for most prizes, but only
    through 1953 for Physiology/Medicine. All nomination-layer edges are
    restricted to these date ranges.

    \item \textbf{Physiology/Medicine vetting body incompleteness.} The
    Swedish Wikipedia source for this committee is self-reportedly
    incomplete.

    \item \textbf{Literature committee extraction fragility.} Committee
    service years were parsed from a Nuxt.js SPA payload, which may
    break if the Swedish Academy redesigns its website.

    \item \textbf{Low nomination QID match rate.} Only 2.3\% of
    nomination archive people were matched to Wikidata, primarily due
    to 68.4\% lacking birth years. This limits cross-layer entity
    resolution in the nomination layer.

    \item \textbf{RSAS end-year imputation.} 19.2\% of RSAS members lack
    end years after death-year imputation. These are treated as currently
    active, which may slightly overcount active members in recent years.

    \item \textbf{Storting term consolidation.} Non-consecutive Storting
    terms are preserved, but gaps between terms may not perfectly reflect
    actual departure and return dates.
\end{enumerate}

\subsection{Assessment of Fitness for Network Analysis}

Despite the structural gaps documented above, the dataset is suitable for
multilayer network analysis of the Nobel Prize selection process for several
reasons. First, the governing body and vetting body layers have high
completeness: QID resolution is 93--100\% across all bodies, and temporal
coverage extends from the founding of each institution to the present (with
the notable exception of Karolinska post-1970). Second, the laureate layer
has 100\% QID coverage and complete temporal span (1901--2025). Third, the
nomination layer, while limited in entity resolution, preserves the full
topological structure of nominator--nominee relationships through the
\texttt{NOM:} identifier system; analyses that do not require cross-layer
linking of nomination-layer individuals can use these edges directly.

The primary analytical limitations are: (a)~Physiology/Medicine analyses
are constrained by the Karolinska gap (no governing body data after 1970)
and the truncated nomination archive (through 1953 only); (b)~cross-layer
analyses linking nominators or nominees to governing/vetting body members
are limited by the 2.3\% QID match rate; and (c)~the 34.2\% isolate node
rate inflates the apparent number of disconnected individuals, though this
is an artifact of the identifier resolution process rather than a true
structural feature of the network. Researchers should consider these
constraints when designing analyses and interpreting results, particularly
for prize-specific comparisons that include Physiology/Medicine.

% =============================================================================

\section{Data Collection and Processing}\label{sec:data}

We constructed a multilayer directed network representing the institutional
structure of the Nobel Prize selection process across all five original prize
categories: Physics, Chemistry, Physiology or Medicine, Literature, and Peace.
The network comprises four layers---governing bodies, vetting bodies
(Nobel Committees), nominators, and nominees/laureates---connected by directed
edges representing institutional authority and participation. All data
gathering was performed programmatically in R, using a reproducible pipeline
of eight scripts. Source code is available at [repository URL]. Below we
describe each stage of the pipeline in sufficient detail for independent
replication.

\subsection{Governing Bodies}\label{sec:governing}

Governing bodies are the institutions that formally award each Nobel Prize.
We gathered membership rosters for four bodies, each mapped to its
corresponding prize category.

\subsubsection{Royal Swedish Academy of Sciences (RSAS)}

The RSAS awards the prizes in Physics and Chemistry. Members are elected for
life; their service begins at induction and ends at death. We scraped the
membership list from the Swedish-language Wikipedia page
\emph{Lista \"over ledam\"oter av Kungliga Vetenskapsakademien}, which
organizes members chronologically under bold year headers (e.g.,
\texttt{<p><b>1740</b></p>}) followed by bulleted lists of member names
linked to individual Wikipedia biography pages. Founding members (1739)
appear under the headings ``Akademiens grundare'' and ``\"Ovriga.''

We parsed the rendered HTML by iterating through child elements of the
main content division in document order, updating the current induction year
whenever a bold year header was encountered, and extracting all hyperlinked
names from subsequent unordered lists. We excluded redlinks (articles not
yet created) and non-article links. For each member, we recorded the
induction year as the start year; end years were left initially blank and
later imputed from Wikidata death dates (Section~\ref{sec:demographics}).

Wikidata QIDs were resolved in batch using the MediaWiki
\texttt{action=query\&prop=pageprops} API, which maps up to 50 Wikipedia
page titles to their corresponding Wikidata items per request. This is
approximately 50 times faster than scraping individual sidebar links. The
API also handles title normalization (e.g., converting underscores to spaces
and resolving Unicode variants), which we accounted for when matching
responses back to input titles. We used exponential backoff (delay doubling
from 2 seconds, up to 5 retries) for failed requests and imposed a
0.2-second delay between batches.

Where multiple records existed for the same QID within the RSAS (e.g., due
to class transfers within the Academy), we retained only the earliest
induction record. The final dataset contains 2{,}312 RSAS members spanning
1739--2025.

\subsubsection{Swedish Academy}

The Swedish Academy awards the Prize in Literature. Members are appointed for
life. We queried Wikidata directly via SPARQL, selecting all individuals
holding position (P39) ``member of the Swedish Academy'' (Q207360), excluding
honorary members (Q97563901). Start and end dates were extracted from the
\texttt{pq:P580} and \texttt{pq:P582} qualifiers on each position statement.
Years were parsed from Wikidata's ISO~8601 datetime format
(e.g., \texttt{1871-12-20T00:00:00Z} $\to$ 1871). The query yielded 202
member records.

All SPARQL queries throughout the pipeline were executed via HTTP POST to the
Wikidata Query Service endpoint (\texttt{https://query.wikidata.org/sparql}),
with results returned in \texttt{application/sparql-results+json} format. We
implemented retry logic with exponential backoff (5-second base delay, up to
10 retries) and identified our bot with a descriptive User-Agent string per
Wikidata's access policy.

\subsubsection{Storting (Norwegian Parliament)}

The Storting awards the Peace Prize. We queried Wikidata for all entities
classified as human (P31 = Q5) holding position (P39) ``member of the
Storting'' (Q9045502), extracting start and end dates from position
qualifiers. Norwegian parliamentarians may serve non-consecutive terms;
we consolidated consecutive terms by grouping each member's service
periods chronologically and merging periods where the gap between one
term's end year and the next term's start year was at most one year.
Non-consecutive terms were preserved as separate records. This yielded
4{,}281 member-term records spanning 1814--2025.

\subsubsection{Karolinska Institutet}

The Karolinska Institutet (KI) awards the Prize in Physiology or Medicine.
Prior to 1977, all full professors at KI constituted the selection body;
in 1977 the Nobel Assembly was established as a distinct entity, and in 1984
its membership was fixed at 50. Neither Wikidata nor KI's website publishes
a complete roster of the Nobel Assembly.

We compiled KI professor data from Project Runeberg's digitized editions of
the Swedish state calendar (\emph{statskalendern}), covering 1881--1970.
The resulting Excel file contains one row per professor with columns for
Wikidata QID, name, and the first and last calendar edition in which each
professor appeared (serving as proxies for start and end years of service).
We retained the 140 records with confirmed QIDs and both start and end years.

\paragraph{Known gap.} Coverage ends at 1970. Post-1970 Nobel Assembly
membership is not publicly available: Wikidata contains essentially no
P39/P463 records for the Nobel Assembly (Q3375124), KI's website does not
publish a member roster, and Project Runeberg's digitized calendars stop
circa 1972. Extending this coverage would require direct contact with the
Nobel Office at KI.

\subsection{Vetting Bodies (Nobel Committees)}\label{sec:vetting}

Each prize has a dedicated Nobel Committee that reviews nominations and
prepares recommendations for the governing body. We gathered membership
data for all five committees.

\subsubsection{Chemistry, Physics, and Physiology/Medicine Committees}

These three committees are documented on Swedish-language Wikipedia. We
extracted membership lists using a general-purpose parser that fetches
wikitext (not rendered HTML) via the MediaWiki API
(\texttt{action=parse\&prop=wikitext}), splits the text into sections by
heading markers (\texttt{==}~and~\texttt{===}), and extracts bullet-list
entries from sections titled ``Tidigare ledam\"oter'' (former members),
``Nuvarande ledam\"oter'' (current members), and
``Sekreterare'' (secretaries).

Each entry follows one of several Swedish Wikipedia conventions:
\texttt{[[Name]], YYYY--YYYY} (former members with date range),
\texttt{[[Name]], invald YYYY} (current members, Swedish format), or
\texttt{[[Name]], YYYY--} (current members, open-ended). We parsed
Wikipedia article titles from double-bracket links (handling both
\texttt{[[Title|Display Name]]} and \texttt{[[Title]]} forms), extracted
start and end years via regex, and resolved QIDs through the batch
MediaWiki API.

The Swedish Wikipedia page for the Physiology/Medicine committee explicitly
states ``Listan \"ar ofullst\"andig'' (the list is incomplete); our data
should be understood in light of this limitation. Final counts: Chemistry 50
members, Physics 51 members, Physiology/Medicine 72 members (59 with QIDs).

\subsubsection{Literature Committee}

The Nobel Committee for Literature is drawn from the Swedish Academy.
No standalone membership list exists on Wikipedia. Instead, we extracted
committee service years from individual member biographies on
\texttt{svenskaakademien.se}.

We first queried Wikidata for all Swedish Academy members with their
seat numbers (extracted from position labels matching the pattern
``member of the Swedish Academy at seat $N$''). We then constructed
biography page URLs following the site's convention:
\texttt{/de-aderton/stol-nr-\{seat\}-\{name-slug\}}, where the name
slug is a lowercased, hyphen-separated, ASCII-transliterated version
of the member's name (e.g., \"a$\to$a, \"o$\to$o, \aa$\to$a).

The Swedish Academy website is built as a Nuxt.js single-page application.
Rather than executing JavaScript, we extracted biography text from the
\texttt{\_\_NUXT\_DATA\_\_} script tag embedded in each page's HTML source.
This payload contains HTML-encoded biography text; we decoded HTML entities,
stripped markup tags, and searched for mentions of ``Nobelkommitt\'en''
followed by year ranges. Two-digit years were expanded to four digits
based on context (e.g., ``1969--86'' $\to$ 1969--1986). We retained only
members whose biographies mentioned Nobel Committee service, yielding 35
members.

\subsubsection{Norwegian Nobel Committee}

The Norwegian Nobel Committee both vets nominations and awards the Peace
Prize. We parsed the membership table from the English Wikipedia page
\emph{List of members of the Norwegian Nobel Committee}. The page contains
a wikitable with columns for member name, start year, end year, tenure,
party affiliation, and chair/deputy status. We extracted the largest table,
parsed start and end years from the second and third columns, and resolved
QIDs from the first hyperlink in each row via the batch MediaWiki API.
This yielded 61 members (57 with QIDs).

\subsection{Nominations}\label{sec:nominations}

We scraped the Nobel Prize Nomination Archive
(\texttt{nobelprize.org/nomination/archive/}), which is subject to a 50-year
secrecy rule. As of 2025, nomination records are available through 1974 or
1975 for Physics, Chemistry, Literature, and Peace. Physiology or Medicine
nominations are available only through 1953; the Karolinska Institutet has
not released records for 1954--1974, likely related to the 1977 Swedish
transparency law and the concurrent restructuring that created the Nobel
Assembly.

The archive's web interface uses numeric prize codes in its URL parameters:
Physics = 1, Chemistry = 2, Physiology or Medicine = 3, Literature = 4,
Peace = 5. For each prize-year combination, we first retrieved the list page
(\texttt{list.php?prize=\{code\}\&year=\{year\}}) and extracted all
nomination IDs from links to detail pages (\texttt{show.php?id=\{id\}}).
We then scraped each detail page to extract the nomination year, prize
category, and the names and internal person IDs of all nominees and
nominators.

Detail pages use a two-column table layout with bold section headers
(``Nominee:'' and ``Nominator:'') followed by person links
(\texttt{show\_people.php?id=\{id\}}). We classified each person link
by walking through all table cells in document order and tracking which
section header was most recently encountered. A single nomination may list
multiple nominees (16.3\% of nominations) and, more rarely, multiple
nominators (1.1\%), reflecting joint nominations and co-signed nomination
letters. For multi-nominee or multi-nominator nominations, we generated
all pairwise nominee--nominator combinations.

Prize names were standardized to a common vocabulary (e.g., ``Physiology
or Medicine'' $\to$ ``Physiology/Medicine''; the Peace Prize archive uses
the header ``Nomination for Nobel Peace Prize'' rather than the
``Nobel Prize in\ldots'' format used by the other four categories, which we
handled as a special case in parsing).

List pages were fetched sequentially with 0.5-second inter-request delays;
detail pages were fetched in parallel across $N-1$ CPU cores (where $N$ is
the number of available cores, capped at 12) with 0.2-second per-worker
delays. Partial results were checkpointed to disk after each prize-year to
allow resumption if the process was interrupted. The final dataset contains
27{,}932 nomination records spanning 22{,}600 unique nominations, 4{,}006
unique nominees, and 11{,}650 unique nominators. Among these, 755 records
(2.7\%) have anonymous or unlisted nominators, and 2 records represent
self-nominations.

\subsection{Laureates}\label{sec:laureates}

We retrieved all Nobel Prize laureates from the official Nobel Prize API
(v2.1, \texttt{api.nobelprize.org/2.1/laureates}). The API provides
structured JSON data including Wikidata QIDs, names, gender, award year,
prize category, prize portion, motivation text, and primary affiliation at
the time of award. We paginated through all results (25 records per page)
and filtered to the five original prize categories, excluding the Sveriges
Riksbank Prize in Economic Sciences (established 1969).

Each laureate record was processed to extract the Wikidata QID (from the
\texttt{wikidata.id} field), full name (preferring \texttt{fullName.en}
over \texttt{knownName.en}), gender, and prize details. Laureates who
received multiple prizes (7 individuals, including Marie Curie, Linus
Pauling, Frederick Sanger, John Bardeen, and K.~Barry Sharpless) appear
as separate records per prize. The final dataset contains 927
laureate-prize records representing 919 unique individuals, with 100\%
QID coverage.

\subsection{Entity Resolution: Matching Nomination Archive People to
Wikidata}\label{sec:entity-resolution}

The nomination archive identifies people by internal numeric IDs and names
only, without Wikidata QIDs or other persistent identifiers. To integrate
nomination data with the rest of the network (which uses Wikidata QIDs as
the canonical identifier), we developed a two-stage matching pipeline.

\subsubsection{Stage 1: Biographical Data Extraction}

For each of the 14{,}471 unique people in the nomination archive, we
scraped their detail page on \texttt{nobelprize.org}
(\texttt{show\_people.php?id=\{id\}}), extracting last name, first name,
gender, birth year, and death year from labeled table cells. Scraping was
parallelized across available cores in chunks of 500, with partial results
checkpointed to disk.

Of the 14{,}471 people, 4{,}571 (31.6\%) had birth years recorded;
the remaining 9{,}903 (68.4\%) lacked birth year data entirely,
rendering them unmatchable through our disambiguation strategy.

Six records contained non-standard gender codes in the source database
(\texttt{<}, \texttt{\S}, \texttt{-}, \texttt{N}, \texttt{J}). We verified
the identities of all six individuals by first name and confirmed they are
male; these were recoded to ``M'' in our output. One record (Antonio
Carrillo Flores) had a transposed death year (1886 instead of 1986); this
was corrected manually.

\subsubsection{Stage 2: Wikidata Matching}

For each person with both a name and birth year (4{,}571 matchable
individuals), we searched the Wikidata API
(\texttt{wbsearchentities}, limited to 10 candidates per query) and then
verified each candidate's birth year via the \texttt{wbgetclaims} API
(property P569, date of birth). A match was accepted only when the
candidate's Wikidata birth year exactly equaled the birth year recorded
in the Nobel archive. We did not attempt fuzzy name matching or
approximate year matching, prioritizing precision over recall.

This procedure matched 339 of 4{,}571 matchable individuals (7.4\%),
representing 2.3\% of all 14{,}474 nomination archive people. The low
match rate reflects several factors: many early-twentieth-century
nominators and nominees lack Wikidata entries; name transliterations
between the archive's conventions and Wikidata's labels may differ; and
our strict birth-year criterion rejects plausible matches where Wikidata
records a different (or no) birth date. Unmatched individuals are
represented in the network using temporary identifiers with a \texttt{NOM:}
prefix (e.g., \texttt{NOM:7283}), preserving the nomination-layer topology
even where cross-layer entity resolution was not possible.

\subsection{Demographic Enrichment}\label{sec:demographics}

We fetched demographic attributes from Wikidata for all individuals with
resolved QIDs across all data sources. The SPARQL query retrieved: English
name label, gender (P21), country of birth (P19, traversed to sovereign
state via P17), nationality (P27), birth date (P569), death date (P570),
occupation (P106), and employer/institution (P108). Queries were executed
in batches of up to 200 QIDs using the \texttt{VALUES} clause, with
individual fallback queries for any batch that failed.

Multi-valued fields (nationality, occupation, institution) were collapsed
into semicolon-delimited strings. Where multiple birth or death dates
existed for a single QID (e.g., due to conflicting sources), we retained
the first non-missing value. Birth and death years were parsed from
Wikidata's ISO~8601 datetime format.

After the initial demographic fetch for governing body, vetting body, and
laureate QIDs, we performed a second pass to fetch demographics for the
339 nomination archive people matched in Section~\ref{sec:entity-resolution},
then re-collapsed all demographic records to produce a single row per QID.

As a post-processing step, we imputed end-of-service years for RSAS members
(who serve for life) using Wikidata death years. Of 2{,}312 RSAS members,
1{,}868 received imputed end years; 444 remain without end years (likely
living members or individuals with unknown death dates).

The final demographics table contains 7{,}221 unique individuals with
coverage rates of 99.8\% for name, 99.5\% for gender (90.3\% male, 9.3\%
female), 78.7\% for birth country, 97.2\% for nationality, 99.5\% for
birth year, 78.2\% for death year, 97.4\% for occupation, and 38.4\% for
institutional affiliation.

\subsection{Network Construction}\label{sec:network}

The final network was assembled from all intermediate data files. Nodes
correspond to unique individuals identified by Wikidata QIDs, with
demographic attributes from Section~\ref{sec:demographics}. Edges are
directed and carry metadata for year, prize category, and the layer
identities of the source and target nodes.

We defined five edge types reflecting the institutional flow of the
selection process:

\begin{enumerate}
    \item \textbf{Governing body $\to$ vetting body.} For each vetting
    body member, we created edges from all governing body members who were
    active at the start of the vetting member's service. ``Active'' means
    the governing body member's start year is at or before the vetting
    member's start year, and the governing body member's end year (if known)
    is at or after the vetting member's start year. The edge year is set to
    the vetting member's start year.

    \item \textbf{Vetting body $\to$ nominator.} For Chemistry, Physics,
    and Physiology/Medicine, Nobel Committees solicit nominations from
    qualified individuals. For each prize-year, we created edges from all
    active vetting body members to all nominators who submitted nominations
    that year. Nominators without resolved QIDs were represented using
    \texttt{NOM:} prefix identifiers.

    \item \textbf{Nominator $\to$ nominee.} Direct edges from each
    nominator to each nominee within a nomination record. For nominations
    with multiple nominees or nominators, all pairwise combinations are
    represented. Edges were deduplicated.

    \item \textbf{Governing body $\to$ laureate.} For Chemistry, Physics,
    Physiology/Medicine, and Literature, we created edges from all governing
    body members active in the award year to the laureate. This represents
    the formal decision-making authority of the governing body.

    \item \textbf{Vetting body $\to$ laureate.} For Peace only, the
    Norwegian Nobel Committee both evaluates nominations and awards the
    prize. Edges connect all active committee members in the award year
    to the laureate.
\end{enumerate}

The body-to-prize mappings are: RSAS $\to$ \{Chemistry, Physics\};
Karolinska Institutet $\to$ Physiology/Medicine; Swedish Academy $\to$
Literature; Storting $\to$ Peace. For vetting bodies: Nobel Committee for
Chemistry $\to$ Chemistry; Nobel Committee for Physics $\to$ Physics; Nobel
Committee for Physiology/Medicine $\to$ Physiology/Medicine; Nobel Committee
for Literature $\to$ Literature; Norwegian Nobel Committee $\to$ Peace.

After deduplication, the network contains 7{,}221 nodes and 305{,}199
edges. Edge counts by type: governing body $\to$ laureate (147{,}965;
48.5\%), vetting body $\to$ nominator (87{,}492; 28.7\%), governing body
$\to$ vetting body (42{,}360; 13.9\%), nominator $\to$ nominee (26{,}654;
8.7\%), and vetting body $\to$ laureate (728; 0.2\%). Edge counts by prize:
Physics (121{,}290), Chemistry (111{,}622), Physiology/Medicine (49{,}060),
Peace (15{,}853), Literature (7{,}374). The lower counts for
Physiology/Medicine reflect the Karolinska data gap post-1970 and the
truncated nomination archive (through 1953 only); the lower count for
Literature reflects the smaller Swedish Academy membership (18 seats).

\subsection{Self-Loops and Dual Roles}\label{sec:selfloops}

The network contains 220 self-loops (edges where the source and target
are the same individual). These arise legitimately from individuals who
occupy roles in multiple layers simultaneously. Specifically: 148
governing $\to$ vetting edges where RSAS members also served on Nobel
Committees; 44 governing $\to$ laureate edges where RSAS members
subsequently won Nobel Prizes (e.g., Svante Arrhenius, Ernest Rutherford);
26 vetting $\to$ nominator edges where committee members also submitted
nominations; and 2 nominator $\to$ nominee self-nominations. We retained
all self-loops as they reflect genuine structural features of the
selection process.

\section{Data Quality Assessment}\label{sec:diagnostics}

We performed a comprehensive 42-point diagnostic audit of all intermediate
and output data files. Below we summarize the results, organized by data
source, and assess the overall fitness of the dataset for network analysis.

\subsection{File Integrity and Completeness}

All nine expected data files (seven intermediate, two final output) were
present and parseable. Row counts, column counts, file sizes, and
modification timestamps were verified. No files were missing or corrupted.

\subsection{Governing Body Quality}

\paragraph{QID completeness.} All four governing bodies achieved 100\% QID
resolution: RSAS (2{,}312/2{,}312), Storting (4{,}281/4{,}281), Swedish
Academy (202/202), and Karolinska Institutet (140/140).

\paragraph{Year coverage.} RSAS membership spans 1739--2025; Storting spans
1814--2025; the Swedish Academy spans 1786--2024; Karolinska Institutet
spans 1884--1970. All year values fall within plausible historical ranges
(1700--2026). No records have start year exceeding end year (after
correction of three RSAS records where death-year imputation produced
anomalous end years; these were reverted to missing).

\paragraph{RSAS end-year imputation.} Of 2{,}312 RSAS members, 1{,}868
received end years imputed from Wikidata death dates. The remaining 444
(19.2\%) lack end years, likely because they are still living or their
death dates are not recorded in Wikidata. These members are treated as
currently active when determining temporal overlap for edge construction.

\paragraph{Multiple terms.} 583 QIDs appear multiple times within the same
governing body. For the Storting, these represent genuinely non-consecutive
parliamentary terms. For RSAS and Swedish Academy, only the earliest record
per QID was retained (duplicates arise from class transfers within the
Academy).

\subsection{Vetting Body Quality}

\paragraph{QID completeness.} Chemistry (50/50, 100\%), Physics (51/51,
100\%), Literature (35/35, 100\%), and the Norwegian Nobel Committee
(57/61, 93.4\%) all achieved high QID resolution. The Physiology/Medicine
committee achieved 59/72 (81.9\%), reflecting the acknowledged
incompleteness of the source list.

\paragraph{Cross-validation with governing bodies.} We verified that
committee members appear in their parent governing body's roster: 42/45
Chemistry committee members appear in RSAS (3 not found), 44/50 Physics
committee members appear in RSAS (6 not found), and 23/52
Physiology/Medicine committee members appear in Karolinska Institutet
(29 not found). The high rate of missing Physiology/Medicine cross-references
is attributable to the Karolinska data covering only 1881--1970, while the
committee list extends to 2025.

\paragraph{Known source limitations.} The Physiology/Medicine committee's
Swedish Wikipedia page states the list is incomplete. The Literature
committee data was extracted from the Swedish Academy's Nuxt.js-based
website by parsing the \texttt{\_\_NUXT\_DATA\_\_} payload, a method that
is fragile to future site redesigns.

\subsection{Nomination Data Quality}

\paragraph{Prize coverage.} All five prize categories are represented:
Physiology/Medicine (6{,}321 records), Peace (6{,}134), Physics (5{,}691),
Chemistry (5{,}094), and Literature (4{,}692). Year ranges align with
expectations: Chemistry, Physics, and Literature cover 1901--1974; Peace
covers 1901--1975; Physiology/Medicine covers 1901--1953.

\paragraph{Missing data.} All 27{,}932 records have nominee person IDs.
755 records (2.7\%) lack nominator person IDs, corresponding to anonymous
or unlisted nominators. Two self-nominations were detected (nominee =
nominator).

\paragraph{Entity resolution.} Of 14{,}474 unique people in the nomination
archive, only 339 (2.3\%) were matched to Wikidata QIDs. The primary
bottleneck is missing biographical data: 9{,}903 people (68.4\%) lack
birth years in the archive, making Wikidata disambiguation infeasible.
Of the 4{,}571 people with sufficient data for matching, 339 (7.4\%) were
successfully resolved. As a consequence, 21.4\% of all edge endpoints
in the final network use unresolved \texttt{NOM:} prefix identifiers, and
35.4\% of edges involve at least one such identifier. These edges preserve
the nomination-layer topology but cannot be linked to individuals in other
layers.

\subsection{Laureate Data Quality}

All 927 laureate-prize records have Wikidata QIDs (100\% coverage), and
all 919 unique laureate QIDs appear in the nodes file. Gender is recorded
for 891 laureates (827 male, 64 female, 28 unknown). Seven individuals
received multiple prizes. All laureates span 1901--2025.

\paragraph{Governing $\to$ laureate edge coverage.} Of 919 laureates, 267
lack governing body $\to$ laureate edges. Of these, 139 are Peace laureates
(expected, since Peace uses the vetting $\to$ laureate pathway), and 128
are Physiology/Medicine laureates awarded after 1970 (the Karolinska data
gap). These 128 missing edge sets correspond to 55 prize-years
(1971--2025) and represent a structural limitation of the available data
rather than a pipeline error.

\subsection{Demographic Data Quality}

The demographics table contains 7{,}221 unique QIDs with no duplicates.
All governing body QIDs (6{,}133), all vetting body QIDs (235), and all
matched nomination QIDs (339) are present. Completeness rates for individual
fields were reported in Section~\ref{sec:demographics}.

\paragraph{Temporal consistency.} No records have birth year exceeding
death year (after correction of one transposed death year). Fifty-seven
records have birth years prior to 1700, corresponding to early members
of the RSAS (founded 1739) and Storting (founded 1814); these are
historically accurate.

\paragraph{Gender distribution.} Demographics report 6{,}518 male (90.3\%),
670 female (9.3\%), and 33 unknown (0.5\%). The nomination archive's own
gender field reports 10{,}054 male (69.5\%), 334 female (2.3\%), and 4{,}086
unknown (28.2\%). The large unknown fraction in the nomination archive
reflects missing biographical data for people without Wikidata entries.

\subsection{Network Structural Quality}

\paragraph{Duplicate edges.} No duplicate edges exist after deduplication.

\paragraph{Isolate nodes.} 2{,}466 of 7{,}221 nodes (34.2\%) have no
edges. These are predominantly nomination archive people who were matched
to Wikidata QIDs (and thus appear in the nodes table) but whose
\texttt{NOM:} prefix identifiers in the edge list were not replaced by
their QIDs. This is a consequence of the network construction pipeline
resolving QIDs for nodes and edges independently; nomination-layer edges
reference the original \texttt{NOM:} identifiers regardless of whether
the person was subsequently matched.

\paragraph{Degree distribution.} Among edges with resolved QIDs on both
endpoints, 4{,}624 nodes are involved. The mean degree is 85.3, the median
is 13, the minimum is 1, and the maximum is 1{,}159 (QID Q4951154). The
high maximum degree reflects governing body members of large institutions
(e.g., RSAS members connected to many laureates and committee members over
long service periods).

\paragraph{Cross-layer entity overlap.} 105 laureates are also governing
body members, 15 laureates served on vetting bodies, 68 laureates appear
as nominees (with resolved QIDs), and 113 individuals appear in multiple
governing bodies. These overlaps reflect the historically intertwined nature
of Swedish scientific institutions and are expected features of the data.

\subsection{Summary of Known Structural Gaps}\label{sec:gaps}

We document seven structural limitations inherent to the available sources,
all of which should be considered when interpreting network analyses:

\begin{enumerate}
    \item \textbf{Karolinska Institutet post-1970 gap.} Nobel Assembly
    membership after 1970 is not publicly available, resulting in missing
    governing body $\to$ laureate and governing body $\to$ vetting body
    edges for Physiology/Medicine after 1970.

    \item \textbf{Nomination archive 50-year secrecy rule.} Nomination
    data are available through 1974--1975 for most prizes, but only
    through 1953 for Physiology/Medicine. All nomination-layer edges are
    restricted to these date ranges.

    \item \textbf{Physiology/Medicine vetting body incompleteness.} The
    Swedish Wikipedia source for this committee is self-reportedly
    incomplete.

    \item \textbf{Literature committee extraction fragility.} Committee
    service years were parsed from a Nuxt.js SPA payload, which may
    break if the Swedish Academy redesigns its website.

    \item \textbf{Low nomination QID match rate.} Only 2.3\% of
    nomination archive people were matched to Wikidata, primarily due
    to 68.4\% lacking birth years. This limits cross-layer entity
    resolution in the nomination layer.

    \item \textbf{RSAS end-year imputation.} 19.2\% of RSAS members lack
    end years after death-year imputation. These are treated as currently
    active, which may slightly overcount active members in recent years.

    \item \textbf{Storting term consolidation.} Non-consecutive Storting
    terms are preserved, but gaps between terms may not perfectly reflect
    actual departure and return dates.
\end{enumerate}

\subsection{Assessment of Fitness for Network Analysis}

Despite the structural gaps documented above, the dataset is suitable for
multilayer network analysis of the Nobel Prize selection process for several
reasons. First, the governing body and vetting body layers have high
completeness: QID resolution is 93--100\% across all bodies, and temporal
coverage extends from the founding of each institution to the present (with
the notable exception of Karolinska post-1970). Second, the laureate layer
has 100\% QID coverage and complete temporal span (1901--2025). Third, the
nomination layer, while limited in entity resolution, preserves the full
topological structure of nominator--nominee relationships through the
\texttt{NOM:} identifier system; analyses that do not require cross-layer
linking of nomination-layer individuals can use these edges directly.

The primary analytical limitations are: (a)~Physiology/Medicine analyses
are constrained by the Karolinska gap (no governing body data after 1970)
and the truncated nomination archive (through 1953 only); (b)~cross-layer
analyses linking nominators or nominees to governing/vetting body members
are limited by the 2.3\% QID match rate; and (c)~the 34.2\% isolate node
rate inflates the apparent number of disconnected individuals, though this
is an artifact of the identifier resolution process rather than a true
structural feature of the network. Researchers should consider these
constraints when designing analyses and interpreting results, particularly
for prize-specific comparisons that include Physiology/Medicine.

% =============================================================================

\section{Data Collection and Processing}\label{sec:data}

We constructed a multilayer directed network representing the institutional
structure of the Nobel Prize selection process across all five original prize
categories: Physics, Chemistry, Physiology or Medicine, Literature, and Peace.
The network comprises four layers---governing bodies, vetting bodies
(Nobel Committees), nominators, and nominees/laureates---connected by directed
edges representing institutional authority and participation. All data
gathering was performed programmatically in R, using a reproducible pipeline
of eight scripts. Source code is available at [repository URL]. Below we
describe each stage of the pipeline in sufficient detail for independent
replication.

\subsection{Governing Bodies}\label{sec:governing}

Governing bodies are the institutions that formally award each Nobel Prize.
We gathered membership rosters for four bodies, each mapped to its
corresponding prize category.

\subsubsection{Royal Swedish Academy of Sciences (RSAS)}

The RSAS awards the prizes in Physics and Chemistry. Members are elected for
life; their service begins at induction and ends at death. We scraped the
membership list from the Swedish-language Wikipedia page
\emph{Lista \"over ledam\"oter av Kungliga Vetenskapsakademien}, which
organizes members chronologically under bold year headers (e.g.,
\texttt{<p><b>1740</b></p>}) followed by bulleted lists of member names
linked to individual Wikipedia biography pages. Founding members (1739)
appear under the headings ``Akademiens grundare'' and ``\"Ovriga.''

We parsed the rendered HTML by iterating through child elements of the
main content division in document order, updating the current induction year
whenever a bold year header was encountered, and extracting all hyperlinked
names from subsequent unordered lists. We excluded redlinks (articles not
yet created) and non-article links. For each member, we recorded the
induction year as the start year; end years were left initially blank and
later imputed from Wikidata death dates (Section~\ref{sec:demographics}).

Wikidata QIDs were resolved in batch using the MediaWiki
\texttt{action=query\&prop=pageprops} API, which maps up to 50 Wikipedia
page titles to their corresponding Wikidata items per request. This is
approximately 50 times faster than scraping individual sidebar links. The
API also handles title normalization (e.g., converting underscores to spaces
and resolving Unicode variants), which we accounted for when matching
responses back to input titles. We used exponential backoff (delay doubling
from 2 seconds, up to 5 retries) for failed requests and imposed a
0.2-second delay between batches.

Where multiple records existed for the same QID within the RSAS (e.g., due
to class transfers within the Academy), we retained only the earliest
induction record. The final dataset contains 2{,}312 RSAS members spanning
1739--2025.

\subsubsection{Swedish Academy}

The Swedish Academy awards the Prize in Literature. Members are appointed for
life. We queried Wikidata directly via SPARQL, selecting all individuals
holding position (P39) ``member of the Swedish Academy'' (Q207360), excluding
honorary members (Q97563901). Start and end dates were extracted from the
\texttt{pq:P580} and \texttt{pq:P582} qualifiers on each position statement.
Years were parsed from Wikidata's ISO~8601 datetime format
(e.g., \texttt{1871-12-20T00:00:00Z} $\to$ 1871). The query yielded 202
member records.

All SPARQL queries throughout the pipeline were executed via HTTP POST to the
Wikidata Query Service endpoint (\texttt{https://query.wikidata.org/sparql}),
with results returned in \texttt{application/sparql-results+json} format. We
implemented retry logic with exponential backoff (5-second base delay, up to
10 retries) and identified our bot with a descriptive User-Agent string per
Wikidata's access policy.

\subsubsection{Storting (Norwegian Parliament)}

The Storting awards the Peace Prize. We queried Wikidata for all entities
classified as human (P31 = Q5) holding position (P39) ``member of the
Storting'' (Q9045502), extracting start and end dates from position
qualifiers. Norwegian parliamentarians may serve non-consecutive terms;
we consolidated consecutive terms by grouping each member's service
periods chronologically and merging periods where the gap between one
term's end year and the next term's start year was at most one year.
Non-consecutive terms were preserved as separate records. This yielded
4{,}281 member-term records spanning 1814--2025.

\subsubsection{Karolinska Institutet}

The Karolinska Institutet (KI) awards the Prize in Physiology or Medicine.
Prior to 1977, all full professors at KI constituted the selection body;
in 1977 the Nobel Assembly was established as a distinct entity, and in 1984
its membership was fixed at 50. Neither Wikidata nor KI's website publishes
a complete roster of the Nobel Assembly.

We compiled KI professor data from Project Runeberg's digitized editions of
the Swedish state calendar (\emph{statskalendern}), covering 1881--1970.
The resulting Excel file contains one row per professor with columns for
Wikidata QID, name, and the first and last calendar edition in which each
professor appeared (serving as proxies for start and end years of service).
We retained the 140 records with confirmed QIDs and both start and end years.

\paragraph{Known gap.} Coverage ends at 1970. Post-1970 Nobel Assembly
membership is not publicly available: Wikidata contains essentially no
P39/P463 records for the Nobel Assembly (Q3375124), KI's website does not
publish a member roster, and Project Runeberg's digitized calendars stop
circa 1972. Extending this coverage would require direct contact with the
Nobel Office at KI.

\subsection{Vetting Bodies (Nobel Committees)}\label{sec:vetting}

Each prize has a dedicated Nobel Committee that reviews nominations and
prepares recommendations for the governing body. We gathered membership
data for all five committees.

\subsubsection{Chemistry, Physics, and Physiology/Medicine Committees}

These three committees are documented on Swedish-language Wikipedia. We
extracted membership lists using a general-purpose parser that fetches
wikitext (not rendered HTML) via the MediaWiki API
(\texttt{action=parse\&prop=wikitext}), splits the text into sections by
heading markers (\texttt{==}~and~\texttt{===}), and extracts bullet-list
entries from sections titled ``Tidigare ledam\"oter'' (former members),
``Nuvarande ledam\"oter'' (current members), and
``Sekreterare'' (secretaries).

Each entry follows one of several Swedish Wikipedia conventions:
\texttt{[[Name]], YYYY--YYYY} (former members with date range),
\texttt{[[Name]], invald YYYY} (current members, Swedish format), or
\texttt{[[Name]], YYYY--} (current members, open-ended). We parsed
Wikipedia article titles from double-bracket links (handling both
\texttt{[[Title|Display Name]]} and \texttt{[[Title]]} forms), extracted
start and end years via regex, and resolved QIDs through the batch
MediaWiki API.

The Swedish Wikipedia page for the Physiology/Medicine committee explicitly
states ``Listan \"ar ofullst\"andig'' (the list is incomplete); our data
should be understood in light of this limitation. Final counts: Chemistry 50
members, Physics 51 members, Physiology/Medicine 72 members (59 with QIDs).

\subsubsection{Literature Committee}

The Nobel Committee for Literature is drawn from the Swedish Academy.
No standalone membership list exists on Wikipedia. Instead, we extracted
committee service years from individual member biographies on
\texttt{svenskaakademien.se}.

We first queried Wikidata for all Swedish Academy members with their
seat numbers (extracted from position labels matching the pattern
``member of the Swedish Academy at seat $N$''). We then constructed
biography page URLs following the site's convention:
\texttt{/de-aderton/stol-nr-\{seat\}-\{name-slug\}}, where the name
slug is a lowercased, hyphen-separated, ASCII-transliterated version
of the member's name (e.g., \"a$\to$a, \"o$\to$o, \aa$\to$a).

The Swedish Academy website is built as a Nuxt.js single-page application.
Rather than executing JavaScript, we extracted biography text from the
\texttt{\_\_NUXT\_DATA\_\_} script tag embedded in each page's HTML source.
This payload contains HTML-encoded biography text; we decoded HTML entities,
stripped markup tags, and searched for mentions of ``Nobelkommitt\'en''
followed by year ranges. Two-digit years were expanded to four digits
based on context (e.g., ``1969--86'' $\to$ 1969--1986). We retained only
members whose biographies mentioned Nobel Committee service, yielding 35
members.

\subsubsection{Norwegian Nobel Committee}

The Norwegian Nobel Committee both vets nominations and awards the Peace
Prize. We parsed the membership table from the English Wikipedia page
\emph{List of members of the Norwegian Nobel Committee}. The page contains
a wikitable with columns for member name, start year, end year, tenure,
party affiliation, and chair/deputy status. We extracted the largest table,
parsed start and end years from the second and third columns, and resolved
QIDs from the first hyperlink in each row via the batch MediaWiki API.
This yielded 61 members (57 with QIDs).

\subsection{Nominations}\label{sec:nominations}

We scraped the Nobel Prize Nomination Archive
(\texttt{nobelprize.org/nomination/archive/}), which is subject to a 50-year
secrecy rule. As of 2025, nomination records are available through 1974 or
1975 for Physics, Chemistry, Literature, and Peace. Physiology or Medicine
nominations are available only through 1953; the Karolinska Institutet has
not released records for 1954--1974, likely related to the 1977 Swedish
transparency law and the concurrent restructuring that created the Nobel
Assembly.

The archive's web interface uses numeric prize codes in its URL parameters:
Physics = 1, Chemistry = 2, Physiology or Medicine = 3, Literature = 4,
Peace = 5. For each prize-year combination, we first retrieved the list page
(\texttt{list.php?prize=\{code\}\&year=\{year\}}) and extracted all
nomination IDs from links to detail pages (\texttt{show.php?id=\{id\}}).
We then scraped each detail page to extract the nomination year, prize
category, and the names and internal person IDs of all nominees and
nominators.

Detail pages use a two-column table layout with bold section headers
(``Nominee:'' and ``Nominator:'') followed by person links
(\texttt{show\_people.php?id=\{id\}}). We classified each person link
by walking through all table cells in document order and tracking which
section header was most recently encountered. A single nomination may list
multiple nominees (16.3\% of nominations) and, more rarely, multiple
nominators (1.1\%), reflecting joint nominations and co-signed nomination
letters. For multi-nominee or multi-nominator nominations, we generated
all pairwise nominee--nominator combinations.

Prize names were standardized to a common vocabulary (e.g., ``Physiology
or Medicine'' $\to$ ``Physiology/Medicine''; the Peace Prize archive uses
the header ``Nomination for Nobel Peace Prize'' rather than the
``Nobel Prize in\ldots'' format used by the other four categories, which we
handled as a special case in parsing).

List pages were fetched sequentially with 0.5-second inter-request delays;
detail pages were fetched in parallel across $N-1$ CPU cores (where $N$ is
the number of available cores, capped at 12) with 0.2-second per-worker
delays. Partial results were checkpointed to disk after each prize-year to
allow resumption if the process was interrupted. The final dataset contains
27{,}932 nomination records spanning 22{,}600 unique nominations, 4{,}006
unique nominees, and 11{,}650 unique nominators. Among these, 755 records
(2.7\%) have anonymous or unlisted nominators, and 2 records represent
self-nominations.

\subsection{Laureates}\label{sec:laureates}

We retrieved all Nobel Prize laureates from the official Nobel Prize API
(v2.1, \texttt{api.nobelprize.org/2.1/laureates}). The API provides
structured JSON data including Wikidata QIDs, names, gender, award year,
prize category, prize portion, motivation text, and primary affiliation at
the time of award. We paginated through all results (25 records per page)
and filtered to the five original prize categories, excluding the Sveriges
Riksbank Prize in Economic Sciences (established 1969).

Each laureate record was processed to extract the Wikidata QID (from the
\texttt{wikidata.id} field), full name (preferring \texttt{fullName.en}
over \texttt{knownName.en}), gender, and prize details. Laureates who
received multiple prizes (7 individuals, including Marie Curie, Linus
Pauling, Frederick Sanger, John Bardeen, and K.~Barry Sharpless) appear
as separate records per prize. The final dataset contains 927
laureate-prize records representing 919 unique individuals, with 100\%
QID coverage.

\subsection{Entity Resolution: Matching Nomination Archive People to
Wikidata}\label{sec:entity-resolution}

The nomination archive identifies people by internal numeric IDs and names
only, without Wikidata QIDs or other persistent identifiers. To integrate
nomination data with the rest of the network (which uses Wikidata QIDs as
the canonical identifier), we developed a two-stage matching pipeline.

\subsubsection{Stage 1: Biographical Data Extraction}

For each of the 14{,}471 unique people in the nomination archive, we
scraped their detail page on \texttt{nobelprize.org}
(\texttt{show\_people.php?id=\{id\}}), extracting last name, first name,
gender, birth year, and death year from labeled table cells. Scraping was
parallelized across available cores in chunks of 500, with partial results
checkpointed to disk.

Of the 14{,}471 people, 4{,}571 (31.6\%) had birth years recorded;
the remaining 9{,}903 (68.4\%) lacked birth year data entirely,
rendering them unmatchable through our disambiguation strategy.

Six records contained non-standard gender codes in the source database
(\texttt{<}, \texttt{\S}, \texttt{-}, \texttt{N}, \texttt{J}). We verified
the identities of all six individuals by first name and confirmed they are
male; these were recoded to ``M'' in our output. One record (Antonio
Carrillo Flores) had a transposed death year (1886 instead of 1986); this
was corrected manually.

\subsubsection{Stage 2: Wikidata Matching}

For each person with both a name and birth year (4{,}571 matchable
individuals), we searched the Wikidata API
(\texttt{wbsearchentities}, limited to 10 candidates per query) and then
verified each candidate's birth year via the \texttt{wbgetclaims} API
(property P569, date of birth). A match was accepted only when the
candidate's Wikidata birth year exactly equaled the birth year recorded
in the Nobel archive. We did not attempt fuzzy name matching or
approximate year matching, prioritizing precision over recall.

This procedure matched 339 of 4{,}571 matchable individuals (7.4\%),
representing 2.3\% of all 14{,}474 nomination archive people. The low
match rate reflects several factors: many early-twentieth-century
nominators and nominees lack Wikidata entries; name transliterations
between the archive's conventions and Wikidata's labels may differ; and
our strict birth-year criterion rejects plausible matches where Wikidata
records a different (or no) birth date. Unmatched individuals are
represented in the network using temporary identifiers with a \texttt{NOM:}
prefix (e.g., \texttt{NOM:7283}), preserving the nomination-layer topology
even where cross-layer entity resolution was not possible.

\subsection{Demographic Enrichment}\label{sec:demographics}

We fetched demographic attributes from Wikidata for all individuals with
resolved QIDs across all data sources. The SPARQL query retrieved: English
name label, gender (P21), country of birth (P19, traversed to sovereign
state via P17), nationality (P27), birth date (P569), death date (P570),
occupation (P106), and employer/institution (P108). Queries were executed
in batches of up to 200 QIDs using the \texttt{VALUES} clause, with
individual fallback queries for any batch that failed.

Multi-valued fields (nationality, occupation, institution) were collapsed
into semicolon-delimited strings. Where multiple birth or death dates
existed for a single QID (e.g., due to conflicting sources), we retained
the first non-missing value. Birth and death years were parsed from
Wikidata's ISO~8601 datetime format.

After the initial demographic fetch for governing body, vetting body, and
laureate QIDs, we performed a second pass to fetch demographics for the
339 nomination archive people matched in Section~\ref{sec:entity-resolution},
then re-collapsed all demographic records to produce a single row per QID.

As a post-processing step, we imputed end-of-service years for RSAS members
(who serve for life) using Wikidata death years. Of 2{,}312 RSAS members,
1{,}868 received imputed end years; 444 remain without end years (likely
living members or individuals with unknown death dates).

The final demographics table contains 7{,}221 unique individuals with
coverage rates of 99.8\% for name, 99.5\% for gender (90.3\% male, 9.3\%
female), 78.7\% for birth country, 97.2\% for nationality, 99.5\% for
birth year, 78.2\% for death year, 97.4\% for occupation, and 38.4\% for
institutional affiliation.

\subsection{Network Construction}\label{sec:network}

The final network was assembled from all intermediate data files. Nodes
correspond to unique individuals identified by Wikidata QIDs, with
demographic attributes from Section~\ref{sec:demographics}. Edges are
directed and carry metadata for year, prize category, and the layer
identities of the source and target nodes.

We defined five edge types reflecting the institutional flow of the
selection process:

\begin{enumerate}
    \item \textbf{Governing body $\to$ vetting body.} For each vetting
    body member, we created edges from all governing body members who were
    active at the start of the vetting member's service. ``Active'' means
    the governing body member's start year is at or before the vetting
    member's start year, and the governing body member's end year (if known)
    is at or after the vetting member's start year. The edge year is set to
    the vetting member's start year.

    \item \textbf{Vetting body $\to$ nominator.} For Chemistry, Physics,
    and Physiology/Medicine, Nobel Committees solicit nominations from
    qualified individuals. For each prize-year, we created edges from all
    active vetting body members to all nominators who submitted nominations
    that year. Nominators without resolved QIDs were represented using
    \texttt{NOM:} prefix identifiers.

    \item \textbf{Nominator $\to$ nominee.} Direct edges from each
    nominator to each nominee within a nomination record. For nominations
    with multiple nominees or nominators, all pairwise combinations are
    represented. Edges were deduplicated.

    \item \textbf{Governing body $\to$ laureate.} For Chemistry, Physics,
    Physiology/Medicine, and Literature, we created edges from all governing
    body members active in the award year to the laureate. This represents
    the formal decision-making authority of the governing body.

    \item \textbf{Vetting body $\to$ laureate.} For Peace only, the
    Norwegian Nobel Committee both evaluates nominations and awards the
    prize. Edges connect all active committee members in the award year
    to the laureate.
\end{enumerate}

The body-to-prize mappings are: RSAS $\to$ \{Chemistry, Physics\};
Karolinska Institutet $\to$ Physiology/Medicine; Swedish Academy $\to$
Literature; Storting $\to$ Peace. For vetting bodies: Nobel Committee for
Chemistry $\to$ Chemistry; Nobel Committee for Physics $\to$ Physics; Nobel
Committee for Physiology/Medicine $\to$ Physiology/Medicine; Nobel Committee
for Literature $\to$ Literature; Norwegian Nobel Committee $\to$ Peace.

After deduplication, the network contains 7{,}221 nodes and 305{,}199
edges. Edge counts by type: governing body $\to$ laureate (147{,}965;
48.5\%), vetting body $\to$ nominator (87{,}492; 28.7\%), governing body
$\to$ vetting body (42{,}360; 13.9\%), nominator $\to$ nominee (26{,}654;
8.7\%), and vetting body $\to$ laureate (728; 0.2\%). Edge counts by prize:
Physics (121{,}290), Chemistry (111{,}622), Physiology/Medicine (49{,}060),
Peace (15{,}853), Literature (7{,}374). The lower counts for
Physiology/Medicine reflect the Karolinska data gap post-1970 and the
truncated nomination archive (through 1953 only); the lower count for
Literature reflects the smaller Swedish Academy membership (18 seats).

\subsection{Self-Loops and Dual Roles}\label{sec:selfloops}

The network contains 220 self-loops (edges where the source and target
are the same individual). These arise legitimately from individuals who
occupy roles in multiple layers simultaneously. Specifically: 148
governing $\to$ vetting edges where RSAS members also served on Nobel
Committees; 44 governing $\to$ laureate edges where RSAS members
subsequently won Nobel Prizes (e.g., Svante Arrhenius, Ernest Rutherford);
26 vetting $\to$ nominator edges where committee members also submitted
nominations; and 2 nominator $\to$ nominee self-nominations. We retained
all self-loops as they reflect genuine structural features of the
selection process.

\section{Data Quality Assessment}\label{sec:diagnostics}

We performed a comprehensive 42-point diagnostic audit of all intermediate
and output data files. Below we summarize the results, organized by data
source, and assess the overall fitness of the dataset for network analysis.

\subsection{File Integrity and Completeness}

All nine expected data files (seven intermediate, two final output) were
present and parseable. Row counts, column counts, file sizes, and
modification timestamps were verified. No files were missing or corrupted.

\subsection{Governing Body Quality}

\paragraph{QID completeness.} All four governing bodies achieved 100\% QID
resolution: RSAS (2{,}312/2{,}312), Storting (4{,}281/4{,}281), Swedish
Academy (202/202), and Karolinska Institutet (140/140).

\paragraph{Year coverage.} RSAS membership spans 1739--2025; Storting spans
1814--2025; the Swedish Academy spans 1786--2024; Karolinska Institutet
spans 1884--1970. All year values fall within plausible historical ranges
(1700--2026). No records have start year exceeding end year (after
correction of three RSAS records where death-year imputation produced
anomalous end years; these were reverted to missing).

\paragraph{RSAS end-year imputation.} Of 2{,}312 RSAS members, 1{,}868
received end years imputed from Wikidata death dates. The remaining 444
(19.2\%) lack end years, likely because they are still living or their
death dates are not recorded in Wikidata. These members are treated as
currently active when determining temporal overlap for edge construction.

\paragraph{Multiple terms.} 583 QIDs appear multiple times within the same
governing body. For the Storting, these represent genuinely non-consecutive
parliamentary terms. For RSAS and Swedish Academy, only the earliest record
per QID was retained (duplicates arise from class transfers within the
Academy).

\subsection{Vetting Body Quality}

\paragraph{QID completeness.} Chemistry (50/50, 100\%), Physics (51/51,
100\%), Literature (35/35, 100\%), and the Norwegian Nobel Committee
(57/61, 93.4\%) all achieved high QID resolution. The Physiology/Medicine
committee achieved 59/72 (81.9\%), reflecting the acknowledged
incompleteness of the source list.

\paragraph{Cross-validation with governing bodies.} We verified that
committee members appear in their parent governing body's roster: 42/45
Chemistry committee members appear in RSAS (3 not found), 44/50 Physics
committee members appear in RSAS (6 not found), and 23/52
Physiology/Medicine committee members appear in Karolinska Institutet
(29 not found). The high rate of missing Physiology/Medicine cross-references
is attributable to the Karolinska data covering only 1881--1970, while the
committee list extends to 2025.

\paragraph{Known source limitations.} The Physiology/Medicine committee's
Swedish Wikipedia page states the list is incomplete. The Literature
committee data was extracted from the Swedish Academy's Nuxt.js-based
website by parsing the \texttt{\_\_NUXT\_DATA\_\_} payload, a method that
is fragile to future site redesigns.

\subsection{Nomination Data Quality}

\paragraph{Prize coverage.} All five prize categories are represented:
Physiology/Medicine (6{,}321 records), Peace (6{,}134), Physics (5{,}691),
Chemistry (5{,}094), and Literature (4{,}692). Year ranges align with
expectations: Chemistry, Physics, and Literature cover 1901--1974; Peace
covers 1901--1975; Physiology/Medicine covers 1901--1953.

\paragraph{Missing data.} All 27{,}932 records have nominee person IDs.
755 records (2.7\%) lack nominator person IDs, corresponding to anonymous
or unlisted nominators. Two self-nominations were detected (nominee =
nominator).

\paragraph{Entity resolution.} Of 14{,}474 unique people in the nomination
archive, only 339 (2.3\%) were matched to Wikidata QIDs. The primary
bottleneck is missing biographical data: 9{,}903 people (68.4\%) lack
birth years in the archive, making Wikidata disambiguation infeasible.
Of the 4{,}571 people with sufficient data for matching, 339 (7.4\%) were
successfully resolved. As a consequence, 21.4\% of all edge endpoints
in the final network use unresolved \texttt{NOM:} prefix identifiers, and
35.4\% of edges involve at least one such identifier. These edges preserve
the nomination-layer topology but cannot be linked to individuals in other
layers.

\subsection{Laureate Data Quality}

All 927 laureate-prize records have Wikidata QIDs (100\% coverage), and
all 919 unique laureate QIDs appear in the nodes file. Gender is recorded
for 891 laureates (827 male, 64 female, 28 unknown). Seven individuals
received multiple prizes. All laureates span 1901--2025.

\paragraph{Governing $\to$ laureate edge coverage.} Of 919 laureates, 267
lack governing body $\to$ laureate edges. Of these, 139 are Peace laureates
(expected, since Peace uses the vetting $\to$ laureate pathway), and 128
are Physiology/Medicine laureates awarded after 1970 (the Karolinska data
gap). These 128 missing edge sets correspond to 55 prize-years
(1971--2025) and represent a structural limitation of the available data
rather than a pipeline error.

\subsection{Demographic Data Quality}

The demographics table contains 7{,}221 unique QIDs with no duplicates.
All governing body QIDs (6{,}133), all vetting body QIDs (235), and all
matched nomination QIDs (339) are present. Completeness rates for individual
fields were reported in Section~\ref{sec:demographics}.

\paragraph{Temporal consistency.} No records have birth year exceeding
death year (after correction of one transposed death year). Fifty-seven
records have birth years prior to 1700, corresponding to early members
of the RSAS (founded 1739) and Storting (founded 1814); these are
historically accurate.

\paragraph{Gender distribution.} Demographics report 6{,}518 male (90.3\%),
670 female (9.3\%), and 33 unknown (0.5\%). The nomination archive's own
gender field reports 10{,}054 male (69.5\%), 334 female (2.3\%), and 4{,}086
unknown (28.2\%). The large unknown fraction in the nomination archive
reflects missing biographical data for people without Wikidata entries.

\subsection{Network Structural Quality}

\paragraph{Duplicate edges.} No duplicate edges exist after deduplication.

\paragraph{Isolate nodes.} 2{,}466 of 7{,}221 nodes (34.2\%) have no
edges. These are predominantly nomination archive people who were matched
to Wikidata QIDs (and thus appear in the nodes table) but whose
\texttt{NOM:} prefix identifiers in the edge list were not replaced by
their QIDs. This is a consequence of the network construction pipeline
resolving QIDs for nodes and edges independently; nomination-layer edges
reference the original \texttt{NOM:} identifiers regardless of whether
the person was subsequently matched.

\paragraph{Degree distribution.} Among edges with resolved QIDs on both
endpoints, 4{,}624 nodes are involved. The mean degree is 85.3, the median
is 13, the minimum is 1, and the maximum is 1{,}159 (QID Q4951154). The
high maximum degree reflects governing body members of large institutions
(e.g., RSAS members connected to many laureates and committee members over
long service periods).

\paragraph{Cross-layer entity overlap.} 105 laureates are also governing
body members, 15 laureates served on vetting bodies, 68 laureates appear
as nominees (with resolved QIDs), and 113 individuals appear in multiple
governing bodies. These overlaps reflect the historically intertwined nature
of Swedish scientific institutions and are expected features of the data.

\subsection{Summary of Known Structural Gaps}\label{sec:gaps}

We document seven structural limitations inherent to the available sources,
all of which should be considered when interpreting network analyses:

\begin{enumerate}
    \item \textbf{Karolinska Institutet post-1970 gap.} Nobel Assembly
    membership after 1970 is not publicly available, resulting in missing
    governing body $\to$ laureate and governing body $\to$ vetting body
    edges for Physiology/Medicine after 1970.

    \item \textbf{Nomination archive 50-year secrecy rule.} Nomination
    data are available through 1974--1975 for most prizes, but only
    through 1953 for Physiology/Medicine. All nomination-layer edges are
    restricted to these date ranges.

    \item \textbf{Physiology/Medicine vetting body incompleteness.} The
    Swedish Wikipedia source for this committee is self-reportedly
    incomplete.

    \item \textbf{Literature committee extraction fragility.} Committee
    service years were parsed from a Nuxt.js SPA payload, which may
    break if the Swedish Academy redesigns its website.

    \item \textbf{Low nomination QID match rate.} Only 2.3\% of
    nomination archive people were matched to Wikidata, primarily due
    to 68.4\% lacking birth years. This limits cross-layer entity
    resolution in the nomination layer.

    \item \textbf{RSAS end-year imputation.} 19.2\% of RSAS members lack
    end years after death-year imputation. These are treated as currently
    active, which may slightly overcount active members in recent years.

    \item \textbf{Storting term consolidation.} Non-consecutive Storting
    terms are preserved, but gaps between terms may not perfectly reflect
    actual departure and return dates.
\end{enumerate}

\subsection{Assessment of Fitness for Network Analysis}

Despite the structural gaps documented above, the dataset is suitable for
multilayer network analysis of the Nobel Prize selection process for several
reasons. First, the governing body and vetting body layers have high
completeness: QID resolution is 93--100\% across all bodies, and temporal
coverage extends from the founding of each institution to the present (with
the notable exception of Karolinska post-1970). Second, the laureate layer
has 100\% QID coverage and complete temporal span (1901--2025). Third, the
nomination layer, while limited in entity resolution, preserves the full
topological structure of nominator--nominee relationships through the
\texttt{NOM:} identifier system; analyses that do not require cross-layer
linking of nomination-layer individuals can use these edges directly.

The primary analytical limitations are: (a)~Physiology/Medicine analyses
are constrained by the Karolinska gap (no governing body data after 1970)
and the truncated nomination archive (through 1953 only); (b)~cross-layer
analyses linking nominators or nominees to governing/vetting body members
are limited by the 2.3\% QID match rate; and (c)~the 34.2\% isolate node
rate inflates the apparent number of disconnected individuals, though this
is an artifact of the identifier resolution process rather than a true
structural feature of the network. Researchers should consider these
constraints when designing analyses and interpreting results, particularly
for prize-specific comparisons that include Physiology/Medicine.

% =============================================================================

\section{Data Collection and Processing}\label{sec:data}

We constructed a multilayer directed network representing the institutional
structure of the Nobel Prize selection process across all five original prize
categories: Physics, Chemistry, Physiology or Medicine, Literature, and Peace.
The network comprises four layers---governing bodies, vetting bodies
(Nobel Committees), nominators, and nominees/laureates---connected by directed
edges representing institutional authority and participation. All data
gathering was performed programmatically in R, using a reproducible pipeline
of eight scripts. Source code is available at [repository URL]. Below we
describe each stage of the pipeline in sufficient detail for independent
replication.

\subsection{Governing Bodies}\label{sec:governing}

Governing bodies are the institutions that formally award each Nobel Prize.
We gathered membership rosters for four bodies, each mapped to its
corresponding prize category.

\subsubsection{Royal Swedish Academy of Sciences (RSAS)}

The RSAS awards the prizes in Physics and Chemistry. Members are elected for
life; their service begins at induction and ends at death. We scraped the
membership list from the Swedish-language Wikipedia page
\emph{Lista \"over ledam\"oter av Kungliga Vetenskapsakademien}, which
organizes members chronologically under bold year headers (e.g.,
\texttt{<p><b>1740</b></p>}) followed by bulleted lists of member names
linked to individual Wikipedia biography pages. Founding members (1739)
appear under the headings ``Akademiens grundare'' and ``\"Ovriga.''

We parsed the rendered HTML by iterating through child elements of the
main content division in document order, updating the current induction year
whenever a bold year header was encountered, and extracting all hyperlinked
names from subsequent unordered lists. We excluded redlinks (articles not
yet created) and non-article links. For each member, we recorded the
induction year as the start year; end years were left initially blank and
later imputed from Wikidata death dates (Section~\ref{sec:demographics}).

Wikidata QIDs were resolved in batch using the MediaWiki
\texttt{action=query\&prop=pageprops} API, which maps up to 50 Wikipedia
page titles to their corresponding Wikidata items per request. This is
approximately 50 times faster than scraping individual sidebar links. The
API also handles title normalization (e.g., converting underscores to spaces
and resolving Unicode variants), which we accounted for when matching
responses back to input titles. We used exponential backoff (delay doubling
from 2 seconds, up to 5 retries) for failed requests and imposed a
0.2-second delay between batches.

Where multiple records existed for the same QID within the RSAS (e.g., due
to class transfers within the Academy), we retained only the earliest
induction record. The final dataset contains 2{,}312 RSAS members spanning
1739--2025.

\subsubsection{Swedish Academy}

The Swedish Academy awards the Prize in Literature. Members are appointed for
life. We queried Wikidata directly via SPARQL, selecting all individuals
holding position (P39) ``member of the Swedish Academy'' (Q207360), excluding
honorary members (Q97563901). Start and end dates were extracted from the
\texttt{pq:P580} and \texttt{pq:P582} qualifiers on each position statement.
Years were parsed from Wikidata's ISO~8601 datetime format
(e.g., \texttt{1871-12-20T00:00:00Z} $\to$ 1871). The query yielded 202
member records.

All SPARQL queries throughout the pipeline were executed via HTTP POST to the
Wikidata Query Service endpoint (\texttt{https://query.wikidata.org/sparql}),
with results returned in \texttt{application/sparql-results+json} format. We
implemented retry logic with exponential backoff (5-second base delay, up to
10 retries) and identified our bot with a descriptive User-Agent string per
Wikidata's access policy.

\subsubsection{Storting (Norwegian Parliament)}

The Storting awards the Peace Prize. We queried Wikidata for all entities
classified as human (P31 = Q5) holding position (P39) ``member of the
Storting'' (Q9045502), extracting start and end dates from position
qualifiers. Norwegian parliamentarians may serve non-consecutive terms;
we consolidated consecutive terms by grouping each member's service
periods chronologically and merging periods where the gap between one
term's end year and the next term's start year was at most one year.
Non-consecutive terms were preserved as separate records. This yielded
4{,}281 member-term records spanning 1814--2025.

\subsubsection{Karolinska Institutet}

The Karolinska Institutet (KI) awards the Prize in Physiology or Medicine.
Prior to 1977, all full professors at KI constituted the selection body;
in 1977 the Nobel Assembly was established as a distinct entity, and in 1984
its membership was fixed at 50. Neither Wikidata nor KI's website publishes
a complete roster of the Nobel Assembly.

We compiled KI professor data from Project Runeberg's digitized editions of
the Swedish state calendar (\emph{statskalendern}), covering 1881--1970.
The resulting Excel file contains one row per professor with columns for
Wikidata QID, name, and the first and last calendar edition in which each
professor appeared (serving as proxies for start and end years of service).
We retained the 140 records with confirmed QIDs and both start and end years.

\paragraph{Known gap.} Coverage ends at 1970. Post-1970 Nobel Assembly
membership is not publicly available: Wikidata contains essentially no
P39/P463 records for the Nobel Assembly (Q3375124), KI's website does not
publish a member roster, and Project Runeberg's digitized calendars stop
circa 1972. Extending this coverage would require direct contact with the
Nobel Office at KI.

\subsection{Vetting Bodies (Nobel Committees)}\label{sec:vetting}

Each prize has a dedicated Nobel Committee that reviews nominations and
prepares recommendations for the governing body. We gathered membership
data for all five committees.

\subsubsection{Chemistry, Physics, and Physiology/Medicine Committees}

These three committees are documented on Swedish-language Wikipedia. We
extracted membership lists using a general-purpose parser that fetches
wikitext (not rendered HTML) via the MediaWiki API
(\texttt{action=parse\&prop=wikitext}), splits the text into sections by
heading markers (\texttt{==}~and~\texttt{===}), and extracts bullet-list
entries from sections titled ``Tidigare ledam\"oter'' (former members),
``Nuvarande ledam\"oter'' (current members), and
``Sekreterare'' (secretaries).

Each entry follows one of several Swedish Wikipedia conventions:
\texttt{[[Name]], YYYY--YYYY} (former members with date range),
\texttt{[[Name]], invald YYYY} (current members, Swedish format), or
\texttt{[[Name]], YYYY--} (current members, open-ended). We parsed
Wikipedia article titles from double-bracket links (handling both
\texttt{[[Title|Display Name]]} and \texttt{[[Title]]} forms), extracted
start and end years via regex, and resolved QIDs through the batch
MediaWiki API.

The Swedish Wikipedia page for the Physiology/Medicine committee explicitly
states ``Listan \"ar ofullst\"andig'' (the list is incomplete); our data
should be understood in light of this limitation. Final counts: Chemistry 50
members, Physics 51 members, Physiology/Medicine 72 members (59 with QIDs).

\subsubsection{Literature Committee}

The Nobel Committee for Literature is drawn from the Swedish Academy.
No standalone membership list exists on Wikipedia. Instead, we extracted
committee service years from individual member biographies on
\texttt{svenskaakademien.se}.

We first queried Wikidata for all Swedish Academy members with their
seat numbers (extracted from position labels matching the pattern
``member of the Swedish Academy at seat $N$''). We then constructed
biography page URLs following the site's convention:
\texttt{/de-aderton/stol-nr-\{seat\}-\{name-slug\}}, where the name
slug is a lowercased, hyphen-separated, ASCII-transliterated version
of the member's name (e.g., \"a$\to$a, \"o$\to$o, \aa$\to$a).

The Swedish Academy website is built as a Nuxt.js single-page application.
Rather than executing JavaScript, we extracted biography text from the
\texttt{\_\_NUXT\_DATA\_\_} script tag embedded in each page's HTML source.
This payload contains HTML-encoded biography text; we decoded HTML entities,
stripped markup tags, and searched for mentions of ``Nobelkommitt\'en''
followed by year ranges. Two-digit years were expanded to four digits
based on context (e.g., ``1969--86'' $\to$ 1969--1986). We retained only
members whose biographies mentioned Nobel Committee service, yielding 35
members.

\subsubsection{Norwegian Nobel Committee}

The Norwegian Nobel Committee both vets nominations and awards the Peace
Prize. We parsed the membership table from the English Wikipedia page
\emph{List of members of the Norwegian Nobel Committee}. The page contains
a wikitable with columns for member name, start year, end year, tenure,
party affiliation, and chair/deputy status. We extracted the largest table,
parsed start and end years from the second and third columns, and resolved
QIDs from the first hyperlink in each row via the batch MediaWiki API.
This yielded 61 members (57 with QIDs).

\subsection{Nominations}\label{sec:nominations}

We scraped the Nobel Prize Nomination Archive
(\texttt{nobelprize.org/nomination/archive/}), which is subject to a 50-year
secrecy rule. As of 2025, nomination records are available through 1974 or
1975 for Physics, Chemistry, Literature, and Peace. Physiology or Medicine
nominations are available only through 1953; the Karolinska Institutet has
not released records for 1954--1974, likely related to the 1977 Swedish
transparency law and the concurrent restructuring that created the Nobel
Assembly.

The archive's web interface uses numeric prize codes in its URL parameters:
Physics = 1, Chemistry = 2, Physiology or Medicine = 3, Literature = 4,
Peace = 5. For each prize-year combination, we first retrieved the list page
(\texttt{list.php?prize=\{code\}\&year=\{year\}}) and extracted all
nomination IDs from links to detail pages (\texttt{show.php?id=\{id\}}).
We then scraped each detail page to extract the nomination year, prize
category, and the names and internal person IDs of all nominees and
nominators.

Detail pages use a two-column table layout with bold section headers
(``Nominee:'' and ``Nominator:'') followed by person links
(\texttt{show\_people.php?id=\{id\}}). We classified each person link
by walking through all table cells in document order and tracking which
section header was most recently encountered. A single nomination may list
multiple nominees (16.3\% of nominations) and, more rarely, multiple
nominators (1.1\%), reflecting joint nominations and co-signed nomination
letters. For multi-nominee or multi-nominator nominations, we generated
all pairwise nominee--nominator combinations.

Prize names were standardized to a common vocabulary (e.g., ``Physiology
or Medicine'' $\to$ ``Physiology/Medicine''; the Peace Prize archive uses
the header ``Nomination for Nobel Peace Prize'' rather than the
``Nobel Prize in\ldots'' format used by the other four categories, which we
handled as a special case in parsing).

List pages were fetched sequentially with 0.5-second inter-request delays;
detail pages were fetched in parallel across $N-1$ CPU cores (where $N$ is
the number of available cores, capped at 12) with 0.2-second per-worker
delays. Partial results were checkpointed to disk after each prize-year to
allow resumption if the process was interrupted. The final dataset contains
27{,}932 nomination records spanning 22{,}600 unique nominations, 4{,}006
unique nominees, and 11{,}650 unique nominators. Among these, 755 records
(2.7\%) have anonymous or unlisted nominators, and 2 records represent
self-nominations.

\subsection{Laureates}\label{sec:laureates}

We retrieved all Nobel Prize laureates from the official Nobel Prize API
(v2.1, \texttt{api.nobelprize.org/2.1/laureates}). The API provides
structured JSON data including Wikidata QIDs, names, gender, award year,
prize category, prize portion, motivation text, and primary affiliation at
the time of award. We paginated through all results (25 records per page)
and filtered to the five original prize categories, excluding the Sveriges
Riksbank Prize in Economic Sciences (established 1969).

Each laureate record was processed to extract the Wikidata QID (from the
\texttt{wikidata.id} field), full name (preferring \texttt{fullName.en}
over \texttt{knownName.en}), gender, and prize details. Laureates who
received multiple prizes (7 individuals, including Marie Curie, Linus
Pauling, Frederick Sanger, John Bardeen, and K.~Barry Sharpless) appear
as separate records per prize. The final dataset contains 927
laureate-prize records representing 919 unique individuals, with 100\%
QID coverage.

\subsection{Entity Resolution: Matching Nomination Archive People to
Wikidata}\label{sec:entity-resolution}

The nomination archive identifies people by internal numeric IDs and names
only, without Wikidata QIDs or other persistent identifiers. To integrate
nomination data with the rest of the network (which uses Wikidata QIDs as
the canonical identifier), we developed a two-stage matching pipeline.

\subsubsection{Stage 1: Biographical Data Extraction}

For each of the 14{,}471 unique people in the nomination archive, we
scraped their detail page on \texttt{nobelprize.org}
(\texttt{show\_people.php?id=\{id\}}), extracting last name, first name,
gender, birth year, and death year from labeled table cells. Scraping was
parallelized across available cores in chunks of 500, with partial results
checkpointed to disk.

Of the 14{,}471 people, 4{,}571 (31.6\%) had birth years recorded;
the remaining 9{,}903 (68.4\%) lacked birth year data entirely,
rendering them unmatchable through our disambiguation strategy.

Six records contained non-standard gender codes in the source database
(\texttt{<}, \texttt{\S}, \texttt{-}, \texttt{N}, \texttt{J}). We verified
the identities of all six individuals by first name and confirmed they are
male; these were recoded to ``M'' in our output. One record (Antonio
Carrillo Flores) had a transposed death year (1886 instead of 1986); this
was corrected manually.

\subsubsection{Stage 2: Wikidata Matching}

For each person with both a name and birth year (4{,}571 matchable
individuals), we searched the Wikidata API
(\texttt{wbsearchentities}, limited to 10 candidates per query) and then
verified each candidate's birth year via the \texttt{wbgetclaims} API
(property P569, date of birth). A match was accepted only when the
candidate's Wikidata birth year exactly equaled the birth year recorded
in the Nobel archive. We did not attempt fuzzy name matching or
approximate year matching, prioritizing precision over recall.

This procedure matched 339 of 4{,}571 matchable individuals (7.4\%),
representing 2.3\% of all 14{,}474 nomination archive people. The low
match rate reflects several factors: many early-twentieth-century
nominators and nominees lack Wikidata entries; name transliterations
between the archive's conventions and Wikidata's labels may differ; and
our strict birth-year criterion rejects plausible matches where Wikidata
records a different (or no) birth date. Unmatched individuals are
represented in the network using temporary identifiers with a \texttt{NOM:}
prefix (e.g., \texttt{NOM:7283}), preserving the nomination-layer topology
even where cross-layer entity resolution was not possible.

\subsection{Demographic Enrichment}\label{sec:demographics}

We fetched demographic attributes from Wikidata for all individuals with
resolved QIDs across all data sources. The SPARQL query retrieved: English
name label, gender (P21), country of birth (P19, traversed to sovereign
state via P17), nationality (P27), birth date (P569), death date (P570),
occupation (P106), and employer/institution (P108). Queries were executed
in batches of up to 200 QIDs using the \texttt{VALUES} clause, with
individual fallback queries for any batch that failed.

Multi-valued fields (nationality, occupation, institution) were collapsed
into semicolon-delimited strings. Where multiple birth or death dates
existed for a single QID (e.g., due to conflicting sources), we retained
the first non-missing value. Birth and death years were parsed from
Wikidata's ISO~8601 datetime format.

After the initial demographic fetch for governing body, vetting body, and
laureate QIDs, we performed a second pass to fetch demographics for the
339 nomination archive people matched in Section~\ref{sec:entity-resolution},
then re-collapsed all demographic records to produce a single row per QID.

As a post-processing step, we imputed end-of-service years for RSAS members
(who serve for life) using Wikidata death years. Of 2{,}312 RSAS members,
1{,}868 received imputed end years; 444 remain without end years (likely
living members or individuals with unknown death dates).

The final demographics table contains 7{,}221 unique individuals with
coverage rates of 99.8\% for name, 99.5\% for gender (90.3\% male, 9.3\%
female), 78.7\% for birth country, 97.2\% for nationality, 99.5\% for
birth year, 78.2\% for death year, 97.4\% for occupation, and 38.4\% for
institutional affiliation.

\subsection{Network Construction}\label{sec:network}

The final network was assembled from all intermediate data files. Nodes
correspond to unique individuals identified by Wikidata QIDs, with
demographic attributes from Section~\ref{sec:demographics}. Edges are
directed and carry metadata for year, prize category, and the layer
identities of the source and target nodes.

We defined five edge types reflecting the institutional flow of the
selection process:

\begin{enumerate}
    \item \textbf{Governing body $\to$ vetting body.} For each vetting
    body member, we created edges from all governing body members who were
    active at the start of the vetting member's service. ``Active'' means
    the governing body member's start year is at or before the vetting
    member's start year, and the governing body member's end year (if known)
    is at or after the vetting member's start year. The edge year is set to
    the vetting member's start year.

    \item \textbf{Vetting body $\to$ nominator.} For Chemistry, Physics,
    and Physiology/Medicine, Nobel Committees solicit nominations from
    qualified individuals. For each prize-year, we created edges from all
    active vetting body members to all nominators who submitted nominations
    that year. Nominators without resolved QIDs were represented using
    \texttt{NOM:} prefix identifiers.

    \item \textbf{Nominator $\to$ nominee.} Direct edges from each
    nominator to each nominee within a nomination record. For nominations
    with multiple nominees or nominators, all pairwise combinations are
    represented. Edges were deduplicated.

    \item \textbf{Governing body $\to$ laureate.} For Chemistry, Physics,
    Physiology/Medicine, and Literature, we created edges from all governing
    body members active in the award year to the laureate. This represents
    the formal decision-making authority of the governing body.

    \item \textbf{Vetting body $\to$ laureate.} For Peace only, the
    Norwegian Nobel Committee both evaluates nominations and awards the
    prize. Edges connect all active committee members in the award year
    to the laureate.
\end{enumerate}

The body-to-prize mappings are: RSAS $\to$ \{Chemistry, Physics\};
Karolinska Institutet $\to$ Physiology/Medicine; Swedish Academy $\to$
Literature; Storting $\to$ Peace. For vetting bodies: Nobel Committee for
Chemistry $\to$ Chemistry; Nobel Committee for Physics $\to$ Physics; Nobel
Committee for Physiology/Medicine $\to$ Physiology/Medicine; Nobel Committee
for Literature $\to$ Literature; Norwegian Nobel Committee $\to$ Peace.

After deduplication, the network contains 7{,}221 nodes and 305{,}199
edges. Edge counts by type: governing body $\to$ laureate (147{,}965;
48.5\%), vetting body $\to$ nominator (87{,}492; 28.7\%), governing body
$\to$ vetting body (42{,}360; 13.9\%), nominator $\to$ nominee (26{,}654;
8.7\%), and vetting body $\to$ laureate (728; 0.2\%). Edge counts by prize:
Physics (121{,}290), Chemistry (111{,}622), Physiology/Medicine (49{,}060),
Peace (15{,}853), Literature (7{,}374). The lower counts for
Physiology/Medicine reflect the Karolinska data gap post-1970 and the
truncated nomination archive (through 1953 only); the lower count for
Literature reflects the smaller Swedish Academy membership (18 seats).

\subsection{Self-Loops and Dual Roles}\label{sec:selfloops}

The network contains 220 self-loops (edges where the source and target
are the same individual). These arise legitimately from individuals who
occupy roles in multiple layers simultaneously. Specifically: 148
governing $\to$ vetting edges where RSAS members also served on Nobel
Committees; 44 governing $\to$ laureate edges where RSAS members
subsequently won Nobel Prizes (e.g., Svante Arrhenius, Ernest Rutherford);
26 vetting $\to$ nominator edges where committee members also submitted
nominations; and 2 nominator $\to$ nominee self-nominations. We retained
all self-loops as they reflect genuine structural features of the
selection process.

\section{Data Quality Assessment}\label{sec:diagnostics}

We performed a comprehensive 42-point diagnostic audit of all intermediate
and output data files. Below we summarize the results, organized by data
source, and assess the overall fitness of the dataset for network analysis.

\subsection{File Integrity and Completeness}

All nine expected data files (seven intermediate, two final output) were
present and parseable. Row counts, column counts, file sizes, and
modification timestamps were verified. No files were missing or corrupted.

\subsection{Governing Body Quality}

\paragraph{QID completeness.} All four governing bodies achieved 100\% QID
resolution: RSAS (2{,}312/2{,}312), Storting (4{,}281/4{,}281), Swedish
Academy (202/202), and Karolinska Institutet (140/140).

\paragraph{Year coverage.} RSAS membership spans 1739--2025; Storting spans
1814--2025; the Swedish Academy spans 1786--2024; Karolinska Institutet
spans 1884--1970. All year values fall within plausible historical ranges
(1700--2026). No records have start year exceeding end year (after
correction of three RSAS records where death-year imputation produced
anomalous end years; these were reverted to missing).

\paragraph{RSAS end-year imputation.} Of 2{,}312 RSAS members, 1{,}868
received end years imputed from Wikidata death dates. The remaining 444
(19.2\%) lack end years, likely because they are still living or their
death dates are not recorded in Wikidata. These members are treated as
currently active when determining temporal overlap for edge construction.

\paragraph{Multiple terms.} 583 QIDs appear multiple times within the same
governing body. For the Storting, these represent genuinely non-consecutive
parliamentary terms. For RSAS and Swedish Academy, only the earliest record
per QID was retained (duplicates arise from class transfers within the
Academy).

\subsection{Vetting Body Quality}

\paragraph{QID completeness.} Chemistry (50/50, 100\%), Physics (51/51,
100\%), Literature (35/35, 100\%), and the Norwegian Nobel Committee
(57/61, 93.4\%) all achieved high QID resolution. The Physiology/Medicine
committee achieved 59/72 (81.9\%), reflecting the acknowledged
incompleteness of the source list.

\paragraph{Cross-validation with governing bodies.} We verified that
committee members appear in their parent governing body's roster: 42/45
Chemistry committee members appear in RSAS (3 not found), 44/50 Physics
committee members appear in RSAS (6 not found), and 23/52
Physiology/Medicine committee members appear in Karolinska Institutet
(29 not found). The high rate of missing Physiology/Medicine cross-references
is attributable to the Karolinska data covering only 1881--1970, while the
committee list extends to 2025.

\paragraph{Known source limitations.} The Physiology/Medicine committee's
Swedish Wikipedia page states the list is incomplete. The Literature
committee data was extracted from the Swedish Academy's Nuxt.js-based
website by parsing the \texttt{\_\_NUXT\_DATA\_\_} payload, a method that
is fragile to future site redesigns.

\subsection{Nomination Data Quality}

\paragraph{Prize coverage.} All five prize categories are represented:
Physiology/Medicine (6{,}321 records), Peace (6{,}134), Physics (5{,}691),
Chemistry (5{,}094), and Literature (4{,}692). Year ranges align with
expectations: Chemistry, Physics, and Literature cover 1901--1974; Peace
covers 1901--1975; Physiology/Medicine covers 1901--1953.

\paragraph{Missing data.} All 27{,}932 records have nominee person IDs.
755 records (2.7\%) lack nominator person IDs, corresponding to anonymous
or unlisted nominators. Two self-nominations were detected (nominee =
nominator).

\paragraph{Entity resolution.} Of 14{,}474 unique people in the nomination
archive, only 339 (2.3\%) were matched to Wikidata QIDs. The primary
bottleneck is missing biographical data: 9{,}903 people (68.4\%) lack
birth years in the archive, making Wikidata disambiguation infeasible.
Of the 4{,}571 people with sufficient data for matching, 339 (7.4\%) were
successfully resolved. As a consequence, 21.4\% of all edge endpoints
in the final network use unresolved \texttt{NOM:} prefix identifiers, and
35.4\% of edges involve at least one such identifier. These edges preserve
the nomination-layer topology but cannot be linked to individuals in other
layers.

\subsection{Laureate Data Quality}

All 927 laureate-prize records have Wikidata QIDs (100\% coverage), and
all 919 unique laureate QIDs appear in the nodes file. Gender is recorded
for 891 laureates (827 male, 64 female, 28 unknown). Seven individuals
received multiple prizes. All laureates span 1901--2025.

\paragraph{Governing $\to$ laureate edge coverage.} Of 919 laureates, 267
lack governing body $\to$ laureate edges. Of these, 139 are Peace laureates
(expected, since Peace uses the vetting $\to$ laureate pathway), and 128
are Physiology/Medicine laureates awarded after 1970 (the Karolinska data
gap). These 128 missing edge sets correspond to 55 prize-years
(1971--2025) and represent a structural limitation of the available data
rather than a pipeline error.

\subsection{Demographic Data Quality}

The demographics table contains 7{,}221 unique QIDs with no duplicates.
All governing body QIDs (6{,}133), all vetting body QIDs (235), and all
matched nomination QIDs (339) are present. Completeness rates for individual
fields were reported in Section~\ref{sec:demographics}.

\paragraph{Temporal consistency.} No records have birth year exceeding
death year (after correction of one transposed death year). Fifty-seven
records have birth years prior to 1700, corresponding to early members
of the RSAS (founded 1739) and Storting (founded 1814); these are
historically accurate.

\paragraph{Gender distribution.} Demographics report 6{,}518 male (90.3\%),
670 female (9.3\%), and 33 unknown (0.5\%). The nomination archive's own
gender field reports 10{,}054 male (69.5\%), 334 female (2.3\%), and 4{,}086
unknown (28.2\%). The large unknown fraction in the nomination archive
reflects missing biographical data for people without Wikidata entries.

\subsection{Network Structural Quality}

\paragraph{Duplicate edges.} No duplicate edges exist after deduplication.

\paragraph{Isolate nodes.} 2{,}466 of 7{,}221 nodes (34.2\%) have no
edges. These are predominantly nomination archive people who were matched
to Wikidata QIDs (and thus appear in the nodes table) but whose
\texttt{NOM:} prefix identifiers in the edge list were not replaced by
their QIDs. This is a consequence of the network construction pipeline
resolving QIDs for nodes and edges independently; nomination-layer edges
reference the original \texttt{NOM:} identifiers regardless of whether
the person was subsequently matched.

\paragraph{Degree distribution.} Among edges with resolved QIDs on both
endpoints, 4{,}624 nodes are involved. The mean degree is 85.3, the median
is 13, the minimum is 1, and the maximum is 1{,}159 (QID Q4951154). The
high maximum degree reflects governing body members of large institutions
(e.g., RSAS members connected to many laureates and committee members over
long service periods).

\paragraph{Cross-layer entity overlap.} 105 laureates are also governing
body members, 15 laureates served on vetting bodies, 68 laureates appear
as nominees (with resolved QIDs), and 113 individuals appear in multiple
governing bodies. These overlaps reflect the historically intertwined nature
of Swedish scientific institutions and are expected features of the data.

\subsection{Summary of Known Structural Gaps}\label{sec:gaps}

We document seven structural limitations inherent to the available sources,
all of which should be considered when interpreting network analyses:

\begin{enumerate}
    \item \textbf{Karolinska Institutet post-1970 gap.} Nobel Assembly
    membership after 1970 is not publicly available, resulting in missing
    governing body $\to$ laureate and governing body $\to$ vetting body
    edges for Physiology/Medicine after 1970.

    \item \textbf{Nomination archive 50-year secrecy rule.} Nomination
    data are available through 1974--1975 for most prizes, but only
    through 1953 for Physiology/Medicine. All nomination-layer edges are
    restricted to these date ranges.

    \item \textbf{Physiology/Medicine vetting body incompleteness.} The
    Swedish Wikipedia source for this committee is self-reportedly
    incomplete.

    \item \textbf{Literature committee extraction fragility.} Committee
    service years were parsed from a Nuxt.js SPA payload, which may
    break if the Swedish Academy redesigns its website.

    \item \textbf{Low nomination QID match rate.} Only 2.3\% of
    nomination archive people were matched to Wikidata, primarily due
    to 68.4\% lacking birth years. This limits cross-layer entity
    resolution in the nomination layer.

    \item \textbf{RSAS end-year imputation.} 19.2\% of RSAS members lack
    end years after death-year imputation. These are treated as currently
    active, which may slightly overcount active members in recent years.

    \item \textbf{Storting term consolidation.} Non-consecutive Storting
    terms are preserved, but gaps between terms may not perfectly reflect
    actual departure and return dates.
\end{enumerate}

\subsection{Assessment of Fitness for Network Analysis}

Despite the structural gaps documented above, the dataset is suitable for
multilayer network analysis of the Nobel Prize selection process for several
reasons. First, the governing body and vetting body layers have high
completeness: QID resolution is 93--100\% across all bodies, and temporal
coverage extends from the founding of each institution to the present (with
the notable exception of Karolinska post-1970). Second, the laureate layer
has 100\% QID coverage and complete temporal span (1901--2025). Third, the
nomination layer, while limited in entity resolution, preserves the full
topological structure of nominator--nominee relationships through the
\texttt{NOM:} identifier system; analyses that do not require cross-layer
linking of nomination-layer individuals can use these edges directly.

The primary analytical limitations are: (a)~Physiology/Medicine analyses
are constrained by the Karolinska gap (no governing body data after 1970)
and the truncated nomination archive (through 1953 only); (b)~cross-layer
analyses linking nominators or nominees to governing/vetting body members
are limited by the 2.3\% QID match rate; and (c)~the 34.2\% isolate node
rate inflates the apparent number of disconnected individuals, though this
is an artifact of the identifier resolution process rather than a true
structural feature of the network. Researchers should consider these
constraints when designing analyses and interpreting results, particularly
for prize-specific comparisons that include Physiology/Medicine.
